% =====================================================================
% Báo cáo BTL MT2013 – Nhóm L16 – XVIII (bản tới Phần 5)
% Biên dịch bằng XeLaTeX, từ thư mục report/latex:
%     xelatex main.tex   (chạy 2 lần để có mục lục)
% Hình được lấy trực tiếp từ results/figures/ (chạy run_all.R trước).
% =====================================================================
\documentclass[12pt,a4paper]{article}

\usepackage[left=2.5cm,right=2cm,top=2.5cm,bottom=2.5cm]{geometry}
\usepackage{amsmath}
\usepackage{newtxmath}        % công thức toán kiểu Times
\usepackage[no-math]{fontspec}
\setmainfont{Times New Roman}
\setsansfont{Arial}
\usepackage{polyglossia}
\setdefaultlanguage{vietnamese}
\newcommand{\coverfont}{}     % bìa dùng chung Times New Roman

\usepackage{icomma}           % dấu phẩy thập phân trong chế độ toán
\usepackage{graphicx}
\usepackage{array,booktabs,multirow,tabularx,longtable,makecell}
\usepackage{caption}
\usepackage{subcaption}
\usepackage{enumitem}
\usepackage{xcolor}
\usepackage{listings}
\usepackage{tikz}
\usepackage{float}
\usepackage[hidelinks]{hyperref}

\newcolumntype{L}[1]{>{\raggedright\arraybackslash}p{#1}}
\graphicspath{{../../results/figures/}{./}}

\captionsetup{labelsep=colon}
\captionsetup[table]{position=above}
\setlength{\parskip}{0.2em}
\renewcommand{\arraystretch}{1.25}
\setlist{itemsep=0.15em}

\addto\captionsvietnamese{%
  \renewcommand{\contentsname}{Mục lục}%
  \renewcommand{\listtablename}{Danh sách bảng}%
  \renewcommand{\listfigurename}{Danh sách hình vẽ}%
  \renewcommand{\tablename}{Bảng}%
  \renewcommand{\figurename}{Hình}%
  \renewcommand{\refname}{TÀI LIỆU THAM KHẢO}%
}

\lstset{
  language=R,
  basicstyle=\ttfamily\small,
  numbers=left, numberstyle=\tiny, numbersep=6pt,
  breaklines=true, columns=fullflexible, keepspaces=true,
  frame=single, rulecolor=\color{gray!60},
  commentstyle=\color{teal!70!black}\itshape,
  keywordstyle=\color{blue!70!black},
  stringstyle=\color{red!60!black},
  extendedchars=true, showstringspaces=false,
  aboveskip=0.8em, belowskip=0.8em
}

% Ký hiệu trung bình và trung vị dùng thống nhất trong bài
\newcommand{\xbar}{\bar{x}}
\newcommand{\xmed}{\tilde{x}}

\begin{document}

% =====================================================================
% TRANG BÌA
% =====================================================================
\begin{titlepage}
\begin{tikzpicture}[remember picture,overlay]
  \draw[line width=1.6pt] ([xshift=1.6cm,yshift=-1.6cm]current page.north west)
        rectangle ([xshift=-1.6cm,yshift=1.6cm]current page.south east);
  \draw[line width=0.5pt] ([xshift=1.75cm,yshift=-1.75cm]current page.north west)
        rectangle ([xshift=-1.75cm,yshift=1.75cm]current page.south east);
\end{tikzpicture}%
\par\vspace*{-\baselineskip}
\coverfont
\setlength{\parindent}{0pt}
\centering
{\large\bfseries ĐẠI HỌC QUỐC GIA THÀNH PHỐ HỒ CHÍ MINH\\[0.3em]
TRƯỜNG ĐẠI HỌC BÁCH KHOA\\[0.3em]
KHOA KHOA HỌC ỨNG DỤNG\par}
\vspace{0.2em}
----------\,$\diamond$\,----------\par
\vspace{1.2cm}
\IfFileExists{logo_bk.png}{\includegraphics[width=3.6cm]{logo_bk.png}}{\rule{0pt}{3.6cm}}\par
\vspace{1.2cm}
{\large\bfseries BÁO CÁO BÀI TẬP LỚN\\[0.3em]
XÁC SUẤT THỐNG KÊ (MT2013)\par}
\vspace{0.6cm}
\rule{0.85\textwidth}{0.6pt}\par
\vspace{0.4cm}
{\large\bfseries PHÂN TÍCH MỐI LIÊN HỆ\\[0.2em]
GIỮA THÀNH PHẦN DINH DƯỠNG, NHÀ SẢN XUẤT\\[0.2em]
VÀ ĐIỂM ĐÁNH GIÁ CỦA NGŨ CỐC ĂN SÁNG (RATING)\par}
\vspace{0.3cm}
\rule{0.85\textwidth}{0.6pt}\par
\vspace{0.6cm}
\begin{tabular}{r l}
\textbf{GVHD:} & TS. Hồ Tuấn Đức \\
\textbf{Nhóm lớp:} & L16 -- Nhóm XVIII -- HK261
\end{tabular}\par
\vspace{0.5cm}
\renewcommand{\arraystretch}{1.1}
\begin{tabular}{|c|l|c|l|}
\hline
\textbf{STT} & \multicolumn{1}{c|}{\textbf{Họ và tên}} & \textbf{MSSV} & \multicolumn{1}{c|}{\textbf{Khoa}} \\
\hline
1 & Lê Chí Khánh Minh       & 2513062 & Quản lý Công nghiệp \\
2 & Nguyễn Phan Thiên Quang & 2513079 & Quản lý Công nghiệp \\
3 & Huỳnh Anh Thư           & 2511984 & Quản lý Công nghiệp \\
4 & Đặng Huỳnh Anh Bảo      & 2410220 & Khoa học Ứng dụng \\
5 & Lê Minh Khánh Duy       & 2512605 & Xây dựng \\
\hline
\end{tabular}\par
\vfill
{\itshape Thành phố Hồ Chí Minh, 2026\par}
\vspace{1cm}
\end{titlepage}

% =====================================================================
% PHẦN ĐẦU
% =====================================================================
\pagenumbering{roman}
\setcounter{page}{1}

\section*{Bảng phân công nhiệm vụ và mức độ đóng góp}
\addcontentsline{toc}{section}{Bảng phân công nhiệm vụ và mức độ đóng góp}

\begin{table}[H]
\centering
\begin{tabular}{|c|L{4.6cm}|c|L{3.6cm}|c|}
\hline
STT & Họ và tên & MSSV & Nhiệm vụ & Mức độ đóng góp \\
\hline
1 & Lê Chí Khánh Minh       & 2513062 & Cơ sở lý thuyết; Tổng hợp & 100\% \\ \hline
2 & Nguyễn Phan Thiên Quang & 2513079 & Cơ sở lý thuyết; Tổng hợp & 100\% \\ \hline
3 & Huỳnh Anh Thư           & 2511984 & Tổng quan dữ liệu; Tổng hợp & 100\% \\ \hline
4 & Đặng Huỳnh Anh Bảo      & 2410220 & Tiền xử lý số liệu; Thống kê mô tả & 100\% \\ \hline
5 & Lê Minh Khánh Duy       & 2512605 & Mục tiêu của BTL; Thống kê mô tả & 100\% \\ \hline
\end{tabular}
\end{table}

\newpage
\section*{Lời cảm ơn}
\addcontentsline{toc}{section}{Lời cảm ơn}

Lời đầu tiên, nhóm chúng em xin gửi lời cảm ơn chân thành và sâu sắc nhất đến thầy
Hồ Tuấn Đức -- giảng viên hướng dẫn môn Xác suất Thống kê. Trong suốt quá trình thực hiện
bài tập lớn, thầy đã tận tình giảng dạy, định hướng đề tài cũng như dành nhiều thời gian góp
ý, sửa chữa các thiếu sót. Những chỉ dẫn quý báu của thầy không chỉ giúp nhóm tháo gỡ các
vướng mắc chuyên môn mà còn hoàn thiện bài báo cáo một cách chỉn chu nhất.

Bên cạnh đó, nhóm cũng xin gửi lời cảm ơn đến quý Thầy Cô trường Đại học Bách khoa -- Đại học
Quốc gia Thành phố Hồ Chí Minh, đặc biệt là các thầy cô thuộc Khoa Khoa học Ứng dụng đã trang
bị cho chúng em nền tảng tri thức và kỹ năng cần thiết trong suốt quá trình học tập. Đây là
cơ sở quan trọng giúp nhóm vận dụng lý thuyết vào việc phân tích, xử lý dữ liệu và hoàn thiện
đề tài này.

Cuối cùng, nhóm xin gửi lời cảm ơn đến các bạn sinh viên đã nhiệt tình chia sẻ tài liệu, trao
đổi kinh nghiệm và hỗ trợ nhóm trong quá trình thu thập thông tin, kiểm tra nội dung. Đồng
thời, sự nỗ lực, tinh thần trách nhiệm và sự phối hợp ăn ý của tất cả các thành viên trong
nhóm đã giúp chúng em không chỉ nâng cao kiến thức chuyên môn mà còn hoàn thiện các kỹ năng
mềm như làm việc nhóm, quản lý thời gian và giải quyết vấn đề.

Dù đã dành nhiều thời gian và nỗ lực để hoàn thành bài tập lớn một cách tốt nhất, song báo
cáo vẫn khó tránh khỏi những hạn chế và thiếu sót. Nhóm rất mong nhận được những ý kiến đóng
góp quý báu từ thầy để đề tài ngày càng hoàn thiện hơn, đồng thời giúp chúng em tích lũy kinh
nghiệm cho các môn học và đề tài nghiên cứu tiếp theo.

Một lần nữa, chúng em xin chân thành cảm ơn!

\begin{flushright}
\textit{Nhóm sinh viên thực hiện đề tài}
\end{flushright}

\newpage
\tableofcontents
\newpage
\addcontentsline{toc}{section}{Danh sách bảng}
\listoftables
\newpage
\addcontentsline{toc}{section}{Danh sách hình vẽ}
\listoffigures
\newpage

\pagenumbering{arabic}
\setcounter{page}{1}

% =====================================================================
\section{TỔNG QUAN DỮ LIỆU}
% =====================================================================

\subsection{Ngữ cảnh, Quy mô và Bản chất Dữ liệu}

Tệp dữ liệu \texttt{cereal.csv} (khai thác từ nguồn \textit{80 Cereals} trên Kaggle
\cite{kaggle}) cung cấp thông tin dinh dưỡng và điểm đánh giá (\texttt{rating}) của các sản
phẩm ngũ cốc ăn sáng thương mại phổ biến tại thị trường Mỹ. Dữ liệu ghi nhận chi tiết về nhà
sản xuất, loại ngũ cốc, các chỉ số dinh dưỡng trên nhãn tiêu dùng, vị trí kệ trưng bày và
điểm đánh giá tổng hợp.

Bộ dữ liệu gốc ở dạng bảng gồm 77 quan trắc (77 sản phẩm từ 7 nhà sản xuất: 74 loại ăn lạnh,
3 loại ăn nóng) và 16 biến thuộc tính. Dữ liệu không chứa ô trống, tuy nhiên xuất hiện các
giá trị mã hóa ẩn bằng số $-1$ đại diện cho dữ liệu khuyết tại 3 biến dinh dưỡng:
\texttt{potass} (2 dòng), \texttt{carbo} (1 dòng) và \texttt{sugars} (1 dòng).

\textbf{Tổng thể và giới hạn của mẫu:} Tổng thể mà nhóm hướng tới là các sản phẩm ngũ cốc ăn
sáng đóng gói bán phổ biến tại Mỹ. Tuy nhiên, 77 sản phẩm trong bộ dữ liệu là một danh sách
các nhãn hàng thông dụng, \emph{không} được chọn ngẫu nhiên từ tổng thể này. Trong khi đó, các
khoảng tin cậy và kiểm định ở phần suy diễn đều giả định mẫu ngẫu nhiên. Vì vậy, các kết luận
suy diễn trong bài chỉ nên được hiểu là áp dụng cho nhóm sản phẩm tương tự với mẫu, và đây là
một giới hạn nhóm ghi nhận ngay từ đầu.

\textbf{Bản chất dữ liệu quan sát:} Đây là tập dữ liệu thứ cấp mang tính chất quan sát, không
phải dữ liệu thu thập qua thiết kế thực nghiệm ngẫu nhiên. Do đó, các phân tích trong nghiên
cứu này chỉ dùng để mô tả và lượng hóa \emph{mối liên hệ} (correlation) giữa các đặc trưng
dinh dưỡng, nhà sản xuất và điểm đánh giá, không suy rộng thành quan hệ nhân quả. Vì lý do
này, tên đề tài dùng cụm từ ``mối liên hệ'' thay cho ``ảnh hưởng''.

\subsection{Phân loại và Vai trò các biến trong nghiên cứu}

Toàn bộ thông tin mô tả, phân loại kiểu biến và vai trò xử lý của 16 biến gốc được tổng hợp
trong Bảng~\ref{tab:bien}.

\begin{table}[H]
\centering
\caption{Từ điển và Vai trò của các biến trong bài tập lớn}
\label{tab:bien}
\small
\begin{tabular}{L{2.3cm} L{1.3cm} L{1.9cm} L{2.2cm} L{5.6cm}}
\toprule
Tên biến & Đơn vị & Kiểu biến gốc & Vai trò phân tích & Phương án xử lý dữ liệu trong BTL \\
\midrule
rating & điểm & Định lượng liên tục & Biến phụ thuộc ($Y$) &
  Giữ nguyên thang đo gốc (18,04 -- 93,70 điểm). \\
\midrule
calories, protein, fat, sodium, fiber, carbo, sugars, potass &
  kcal, g, mg & Định lượng liên tục & Biến độc lập định lượng ($X$) &
  Chuẩn hóa hàm lượng về đơn vị 1 oz ($X_{\text{norm}} = X/\text{weight}$) để so sánh
  trên cùng quy mô khẩu phần. \\
\midrule
vitamins & \% & Định lượng rời rạc & Biến độc lập định lượng ($X$) &
  Nhận 3 giá trị (0, 25, 100\%); giữ nguyên, không chia cho weight (xem mục~\ref{sec:canc}). \\
\midrule
mfr & -- & Định tính định danh & Biến phân nhóm &
  So sánh giữa 6 hãng ($N = 76$). Lọc bỏ hãng A do $n_A = 1$. \\
\midrule
shelf & -- & Định tính thứ bậc & Biến phân nhóm &
  So sánh giữa 3 mức vị trí kệ (1, 2, 3). \\
\midrule
weight, cups & oz, cup & Định lượng liên tục & Quy đổi quy mô &
  weight dùng để chia chuẩn hóa dinh dưỡng về 1 oz, sau đó bỏ. \\
\midrule
name, type & -- & Định danh / Nhị phân & Bỏ / Không phân nhóm &
  name chỉ định danh; type chỉ có 3/77 mẫu ăn nóng (3,9\%, quá lệch mẫu). \\
\bottomrule
\end{tabular}
\end{table}

\subsection{Căn cứ Phương pháp luận Xử lý và Chuyển đổi Biến}
\label{sec:canc}

Dựa trên đặc điểm cấu trúc dữ liệu ở Bảng~\ref{tab:bien}, nhóm đề xuất 4 nguyên tắc xử lý
phương pháp luận cho các bước tiếp theo:

\begin{enumerate}
  \item \textbf{Chuẩn hóa quy mô khẩu phần (1 oz):} Khối lượng một khẩu phần gốc
  (\texttt{weight}) giữa các loại ngũ cốc biến thiên từ 0,5 oz đến 1,5 oz (chênh lệch gấp 3
  lần). Việc chia cho \texttt{weight} giúp 8 chỉ số dinh dưỡng được đưa về hàm lượng trên 1 oz,
  đảm bảo mọi sản phẩm được so sánh trên cùng một quy mô. Riêng \texttt{vitamins} là mức
  \% nhu cầu khuyến nghị được ghi theo bậc cố định (0, 25, 100\%) trên nhãn, không phải một
  hàm lượng đo liên tục, nên nhóm giữ nguyên thay vì chia cho \texttt{weight}.

  \item \textbf{Kiểm soát đa cộng tuyến hoàn hảo:} Biến khám phá
  \texttt{total\_carbs} (= \texttt{carbo} + \texttt{sugars}) chỉ phục vụ Thống kê mô tả;
  không đưa đồng thời cả 3 biến (\texttt{carbo}, \texttt{sugars}, \texttt{total\_carbs}) vào
  cùng một mô hình hồi quy vì ma trận $\mathbf{X}^\top\mathbf{X}$ sẽ suy biến.

  \item \textbf{Lọc mẫu hãng A ($n_A = 1$) và xử lý biến type:} Nhà sản xuất A chỉ có duy
  nhất 1 sản phẩm (Maypo). ANOVA vẫn tính được với nhóm một phần tử (nhóm này góp 0 bậc tự do
  vào $SSW$), nhưng không thể kiểm tra giả định chuẩn và phương sai cho nhóm đó, cũng không
  thể so sánh cặp. Phần tương quan và hồi quy về nguyên tắc có thể dùng đủ $N = 77$; tuy
  nhiên nhóm chọn dùng \emph{một mẫu phân tích chung} $N = 76$ cho mọi phần để các kết quả mô
  tả, so sánh nhóm và hồi quy so sánh được với nhau. Biến \texttt{type} chỉ có 3/77 mẫu ăn
  nóng (3,9\%), quá mất cân bằng nên không dùng làm biến phân nhóm.

  \item \textbf{Xử lý biến vitamins và đa cộng tuyến dinh dưỡng:} Biến \texttt{vitamins}
  được giữ định dạng định lượng rời rạc (0, 25, 100\%) để đưa trực tiếp vào hồi quy. Ngoài
  ra, do năng lượng (\texttt{calories}) được cấu thành từ protein, fat và carbo, nhóm sẽ đánh
  giá chỉ số phóng đại phương sai (VIF) để chẩn đoán đa cộng tuyến trước khi chốt mô hình.
\end{enumerate}

\subsection{Lưu ý về bản chất của biến rating}
\label{sec:rating}

Trong quá trình kiểm tra dữ liệu, nhóm hồi quy \texttt{rating} theo 9 biến dinh dưỡng
\emph{gốc} (tính theo khẩu phần, chưa chia \texttt{weight}) trên cả 77 quan trắc và thu
được $R^2 = 1{,}000000$, phần dư lớn nhất chỉ khoảng $5{,}7\cdot10^{-7}$ (sai số làm tròn
ở chữ số thập phân thứ 6 của rating):
\begin{multline}
\text{rating} = 54{,}93 - 0{,}223\,\text{calories} + 3{,}273\,\text{protein}
  - 1{,}691\,\text{fat} - 0{,}0545\,\text{sodium} \\
  + 3{,}444\,\text{fiber} + 1{,}093\,\text{carbo} - 0{,}725\,\text{sugars}
  - 0{,}034\,\text{potass} - 0{,}051\,\text{vitamins}
\end{multline}
Công thức này còn được tính với chính các mã khuyết $-1$ như số thật. Như vậy
\texttt{rating} \emph{không} phải là một điểm đánh giá độc lập, mà là một điểm tính sẵn từ
các cột dinh dưỡng theo khẩu phần. Phát hiện này có ba hệ quả nhóm ghi nhận cho các phần sau.

Hồi quy trên dữ liệu gốc không có sai số ngẫu nhiên, nên các kiểm định $t$, $F$ và
phần chẩn đoán phần dư ở mục~\ref{sec:ols} sẽ không còn ý nghĩa trên dữ liệu đó.

\texttt{mfr} và \texttt{shelf} không có tác động riêng lên \texttt{rating}; mọi chênh
lệch rating giữa các hãng và các kệ phản ánh khác biệt về thành phần dinh dưỡng.

Vì rating tính theo khẩu phần, sau khi chuẩn hóa về 1 oz và điền khuyết, cùng mô hình
chỉ còn $R^2 = 0{,}979$; phần dư của mô hình trên dữ liệu 1 oz đến từ bước chuẩn hóa (13 sản
phẩm có weight $\neq 1$) và điền khuyết, không phải sai số đo.
Phần Thống kê mô tả (Phần~\ref{sec:mota}) không bị ảnh hưởng. Hướng xử lý cụ thể cho phần
suy diễn sẽ được nhóm thống nhất với giảng viên trước khi thực hiện.

% =====================================================================
\section{MỤC TIÊU CỦA BÀI TẬP LỚN}
% =====================================================================

\subsection{Mục tiêu tổng quát}

Nghiên cứu ứng dụng các phương pháp Thống kê Mô tả và Thống kê Suy diễn để khám phá mối liên
hệ giữa các thành phần dinh dưỡng đã chuẩn hóa theo 1 oz với điểm đánh giá sản phẩm
(\texttt{rating}), đồng thời kiểm định và so sánh sự khác biệt về điểm đánh giá và hàm lượng
dinh dưỡng chuẩn hóa giữa các nhóm nhà sản xuất (\texttt{mfr}) và vị trí kệ trưng bày
(\texttt{shelf}).

\subsection{Các mục tiêu cụ thể, hoạt động thực hiện và phương pháp}

\textbf{A. Thống kê Mô tả và Tiền xử lý Dữ liệu}
\begin{itemize}
  \item \textbf{Quy trình tiền xử lý dữ liệu:}
  \begin{enumerate}
    \item \textit{Lọc mẫu sơ bộ:} Loại bỏ quan trắc của nhà sản xuất A ($n_A = 1$) vì không
    kiểm tra được giả định và không so sánh cặp được cho nhóm này, giữ lại mẫu phân tích chung
    gồm $N = 76$ quan trắc.
    \item \textit{Điền khuyết dữ liệu:} Nhận diện và xử lý các giá trị khuyết mã ẩn $-1$ tại
    các biến \texttt{carbo}, \texttt{sugars}, \texttt{potass} bằng trung vị tính trên đúng mẫu
    phân tích $N = 76$.
    \item \textit{Chuẩn hóa quy mô khẩu phần:} Quy đổi 8 biến dinh dưỡng về đơn vị trên 1 oz
    ($X_{\text{norm}} = X/\text{weight}$); sau đó loại bỏ \texttt{weight}, \texttt{cups},
    \texttt{name}, \texttt{type} và tạo biến tổng \texttt{total\_carbs} phục vụ Thống kê mô tả.
  \end{enumerate}
  \item \textbf{Mô tả đặc tính dữ liệu:} Tính các đại lượng đo lường xu hướng trung tâm
  (Trung bình, Trung vị) và độ biến động (Độ lệch chuẩn, IQR, hệ số bất đối xứng), kết hợp
  trực quan hóa phân bố cho biến phụ thuộc \texttt{rating} và các biến dinh dưỡng.
  \begin{itemize}
    \item \textit{Phương pháp:} Thống kê mô tả, biểu đồ phân bố (Histogram + KDE), biểu đồ hộp
    (Boxplot), ma trận tương quan Pearson và Spearman.
  \end{itemize}
\end{itemize}

\textbf{B. Phân tích So sánh Nhóm (Thống kê Suy diễn)}
\begin{itemize}
  \item \textbf{So sánh theo Nhà sản xuất (mfr -- 6 nhóm):} Kiểm định sự khác biệt về điểm
  \texttt{rating} giữa 6 nhà sản xuất.
  \begin{itemize}
    \item \textit{Phương pháp:} Kiểm định tính chuẩn (Shapiro-Wilk) trên từng nhóm và tính
    đồng nhất phương sai (Levene, biến thể Brown-Forsythe). Nếu thỏa mãn giả định, áp dụng
    One-way ANOVA và so sánh bội Tukey HSD; nếu vi phạm giả định chuẩn, dùng Kruskal-Wallis kèm
    so sánh hậu nghiệm Dunn (hoặc Welch's ANOVA nếu phương sai không đồng nhất).
  \end{itemize}
  \item \textbf{So sánh theo Vị trí kệ trưng bày (shelf -- 3 nhóm):} Kiểm định sự khác biệt
  về \texttt{rating} giữa 3 mức kệ (Kệ 1: thấp, Kệ 2: vừa, Kệ 3: cao).
  \begin{itemize}
    \item \textit{Phương pháp:} One-way ANOVA hoặc Kruskal-Wallis tương ứng với mức độ thỏa
    mãn các giả định.
  \end{itemize}
  \item \textbf{Kiểm soát kiểm định nhiều lần:} Nếu mở rộng so sánh sang 8 chất dinh dưỡng
  theo cả \texttt{mfr} lẫn \texttt{shelf} (khoảng 18 kiểm định), xác suất có ít nhất một kết
  luận sai ở $\alpha = 0{,}05$ sẽ cộng dồn đáng kể. Nhóm sẽ hiệu chỉnh $p$-value bằng phương
  pháp Holm (hoặc Bonferroni), hoặc giới hạn số biến đem kiểm định.
  \item \textbf{Diễn giải:} Theo mục~\ref{sec:rating}, khác biệt rating giữa các nhóm (nếu
  có) được diễn giải là khác biệt về thành phần dinh dưỡng giữa các hãng/kệ, không phải tác
  động riêng của thương hiệu hay vị trí kệ.
\end{itemize}

\textbf{C. Phân tích Mối liên hệ và Mô hình hóa (Tương quan \& Hồi quy)}
\begin{itemize}
  \item \textbf{Phân tích tương quan \& Kiểm tra Đa cộng tuyến:} Đánh giá mức độ và chiều
  hướng tương quan giữa từng thành phần dinh dưỡng với \texttt{rating}; kiểm tra hiện tượng
  đa cộng tuyến giữa các biến dinh dưỡng.
  \begin{itemize}
    \item \textit{Phương pháp:} Hệ số tương quan Pearson và Spearman, kiểm định ý nghĩa hệ số
    tương quan; chỉ số phóng đại phương sai (VIF).
  \end{itemize}
  \item \textbf{Mô hình Hồi quy Tuyến tính Bội:} Xây dựng mô hình hồi quy để lượng hóa trọng
  số của từng thành phần dinh dưỡng trong \texttt{rating}.
  \begin{itemize}
    \item \textit{Phương pháp:} Bình phương bé nhất (OLS), $R^2$, $R^2_{\text{adj}}$, kiểm định
    $t$, $F$ và phân tích phần dư. Phạm vi suy diễn trên hệ số sẽ được điều chỉnh theo
    mục~\ref{sec:rating}.
  \end{itemize}
\end{itemize}

% =====================================================================
\section{CƠ SỞ LÝ THUYẾT}
% =====================================================================

\subsection{Thống kê mô tả}

Thống kê mô tả cung cấp các công cụ toán học và đồ họa để tóm tắt các đặc tính cơ bản của tập
dữ liệu mẫu trước khi thực hiện các phân tích suy diễn. Trong toàn bài, $\xbar$ ký hiệu trung
bình mẫu và $\xmed$ ký hiệu trung vị mẫu.

\subsubsection{Các đại lượng đo lường xu hướng trung tâm}

\begin{itemize}
  \item \textbf{Trung bình mẫu (Sample Mean, $\xbar$):} Đại diện cho trọng tâm của tập dữ liệu:
  \begin{equation}
    \xbar = \frac{1}{n}\sum_{i=1}^{n} x_i
  \end{equation}
  \textit{Phạm vi áp dụng:} Biến định lượng có phân phối tương đối đối xứng và ít bị ảnh hưởng
  bởi giá trị ngoại lai.

  \item \textbf{Trung vị mẫu (Sample Median, $\xmed$):} Giá trị nằm ở chính giữa chuỗi dữ liệu
  đã sắp xếp tăng dần $x_{(1)} \le x_{(2)} \le \dots \le x_{(n)}$:
  \begin{equation}
    \xmed =
    \begin{cases}
      x_{\left(\frac{n+1}{2}\right)} & \text{khi } n \text{ lẻ}\\[0.4em]
      \dfrac{1}{2}\left(x_{\left(\frac{n}{2}\right)} + x_{\left(\frac{n}{2}+1\right)}\right) & \text{khi } n \text{ chẵn}
    \end{cases}
  \end{equation}
  \textit{Phạm vi áp dụng:} Thước đo vững (robust), thích hợp khi dữ liệu có phân phối lệch
  hoặc chứa điểm dị biệt.

  \item \textbf{Yếu vị (Mode, $M_0$):} Giá trị xuất hiện với tần số lớn nhất trong tập quan
  trắc, áp dụng cho cả biến định lượng rời rạc và biến định tính.
\end{itemize}

\subsubsection{Các đại lượng đo lường độ phân tán và hình dạng phân phối}

\begin{itemize}
  \item \textbf{Phương sai mẫu hiệu chỉnh ($s^2$):}
  \begin{equation}
    s^2 = \frac{1}{n-1}\sum_{i=1}^{n}(x_i - \xbar)^2
  \end{equation}
  \item \textbf{Độ lệch chuẩn mẫu ($s$):} Căn bậc hai của phương sai mẫu, cùng đơn vị đo với
  biến gốc:
  \begin{equation}
    s = \sqrt{s^2}
  \end{equation}
  \item \textbf{Khoảng biến thiên (Range, $R$):} $R = x_{\max} - x_{\min}$.
  \item \textbf{Tứ phân vị và Khoảng tứ phân vị (IQR):} Tứ phân vị thứ nhất $Q_1$ tại vị trí
  25\%, tứ phân vị thứ ba $Q_3$ tại vị trí 75\%. Nhóm tính phân vị mức $p$ theo phương pháp
  nội suy tuyến tính:
  \begin{equation}
    Q(p) = x_{(\lfloor h \rfloor)} + (h - \lfloor h \rfloor)\left(x_{(\lfloor h \rfloor + 1)} - x_{(\lfloor h \rfloor)}\right),
    \qquad h = (n-1)p + 1
  \end{equation}
  \begin{equation}
    IQR = Q_3 - Q_1
  \end{equation}
  $IQR$ thể hiện độ biến thiên của 50\% dữ liệu ở phần trung tâm, không bị ảnh hưởng bởi các
  giá trị ngoại lai. Lưu ý: cách tính phân vị khác (ví dụ trong một số giáo trình) có thể cho
  $Q_1$, $Q_3$ hơi khác, kéo theo hàng rào Tukey ở mục~\ref{sec:tukey} thay đổi.

  \item \textbf{Hệ số bất đối xứng (Skewness, $G_1$):} Đo mức độ lệch của phân phối. Nhóm dùng
  hệ số có hiệu chỉnh mẫu:
  \begin{equation}
    g_1 = \frac{\frac{1}{n}\sum_{i=1}^{n}(x_i-\xbar)^3}{\left[\frac{1}{n}\sum_{i=1}^{n}(x_i-\xbar)^2\right]^{3/2}},
    \qquad
    G_1 = \frac{\sqrt{n(n-1)}}{n-2}\, g_1
  \end{equation}
  Quy ước phân loại dùng trong bài: $|G_1| < 0{,}5$: tương đối đối xứng;
  $0{,}5 \le |G_1| \le 1$: lệch vừa; $|G_1| > 1$: lệch mạnh. $G_1 > 0$ là lệch phải, $G_1 < 0$
  là lệch trái.
\end{itemize}

\subsubsection{Phương pháp phát hiện giá trị ngoại lai theo hàng rào Tukey}
\label{sec:tukey}

Một quan trắc $x_i$ được xác định là điểm dị biệt nếu nằm ngoài khoảng giá trị hợp lệ:
\begin{equation}
  \left[\,Q_1 - 1{,}5 \times IQR;\; Q_3 + 1{,}5 \times IQR\,\right]
\end{equation}
\begin{itemize}
  \item Điểm ngoại lai dưới: $x_i < Q_1 - 1{,}5 \times IQR$.
  \item Điểm ngoại lai trên: $x_i > Q_3 + 1{,}5 \times IQR$.
\end{itemize}

\subsubsection{Thống kê mô tả cho biến định tính}

Đối với các biến định tính (\texttt{mfr}, \texttt{shelf}), phân tích mô tả sử dụng bảng tần số:
\begin{itemize}
  \item \textbf{Tần số tuyệt đối ($n_k$):} Số lần xuất hiện của nhóm $k$.
  \item \textbf{Tần số tương đối / Tỷ lệ phần trăm ($f_k$):}
  \begin{equation}
    f_k = \frac{n_k}{N}\times 100\%, \quad \text{với } N = \sum_{k=1}^{K} n_k
    \text{ và } \sum_{k=1}^{K} f_k = 100\%
  \end{equation}
\end{itemize}

\subsection{Lý thuyết Kiểm định giả thuyết thống kê \& Ước lượng}

\subsubsection{Nguyên lý kiểm định giả thuyết}
\begin{itemize}
  \item \textbf{Giả thuyết không ($H_0$):} Phát biểu mặc định rằng không có sự khác biệt hoặc
  không có liên hệ.
  \item \textbf{Đối thuyết ($H_1$):} Phát biểu phản bác lại $H_0$, đại diện cho giả định
  nghiên cứu.
  \item \textbf{Sai lầm trong kiểm định:}
  \begin{itemize}
    \item Sai lầm loại I ($\alpha$): Bác bỏ $H_0$ khi $H_0$ đúng:
    $\alpha = P(\text{Bác bỏ } H_0 \mid H_0 \text{ đúng})$.
    \item Sai lầm loại II ($\beta$): Không bác bỏ $H_0$ khi $H_0$ sai:
    $\beta = P(\text{Không bác bỏ } H_0 \mid H_0 \text{ sai})$.
  \end{itemize}
  \item \textbf{Tiêu chuẩn quyết định dựa vào $p$-value:}
  \begin{itemize}
    \item $p\text{-value} < \alpha$: Bác bỏ $H_0$, kết quả có ý nghĩa thống kê.
    \item $p\text{-value} \ge \alpha$: Chưa đủ bằng chứng bác bỏ $H_0$.
  \end{itemize}
  \item \textbf{Kiểm định nhiều lần:} Khi thực hiện $m$ kiểm định cùng lúc, xác suất mắc ít
  nhất một sai lầm loại I tăng lên (với $m$ kiểm định độc lập là $1-(1-\alpha)^m$). Phương pháp
  Holm sắp các $p$-value tăng dần $p_{(1)} \le \dots \le p_{(m)}$ và bác bỏ lần lượt
  $H_{(j)}$ khi $p_{(j)} \le \alpha/(m - j + 1)$, dừng ở giả thuyết đầu tiên không bác bỏ được.
\end{itemize}

\subsubsection{Ước lượng khoảng tin cậy}

Bảng~\ref{tab:ktc} tổng hợp các bài toán tìm khoảng tin cậy đối xứng độ tin cậy
$\gamma = 1-\alpha$ cho các tham số đặc trưng của tổng thể.

\begin{table}[H]
\centering
\caption{Tóm tắt các bài toán tìm khoảng tin cậy đối xứng}
\label{tab:ktc}
\small
\begin{tabular}{L{3cm} L{4.2cm} L{3.4cm} L{3.8cm}}
\toprule
Dạng tham số & Giả định mẫu / tổng thể & Độ chính xác ($\varepsilon$) & Khoảng tin cậy \\
\midrule
Tỷ lệ ($p$) & $nf \ge 5$ và $n(1-f) \ge 5$ &
  $z_{\alpha/2}\sqrt{\dfrac{f(1-f)}{n}}$ & $f - \varepsilon < p < f + \varepsilon$ \\
\midrule
Trung bình ($\mu$) -- TH1 & $X \sim N(\mu,\sigma^2)$, $\sigma$ đã biết &
  $z_{\alpha/2}\dfrac{\sigma}{\sqrt{n}}$ & $\xbar - \varepsilon < \mu < \xbar + \varepsilon$ \\
\midrule
Trung bình ($\mu$) -- TH2 & $X \sim N(\mu,\sigma^2)$, $\sigma$ chưa biết &
  $t_{\alpha/2;\,n-1}\dfrac{s}{\sqrt{n}}$ & $\xbar - \varepsilon < \mu < \xbar + \varepsilon$ \\
\midrule
Trung bình ($\mu$) -- TH3 & Phân phối tùy ý, mẫu lớn $n \ge 30$ &
  $z_{\alpha/2}\dfrac{s}{\sqrt{n}}$ & $\xbar - \varepsilon < \mu < \xbar + \varepsilon$ \\
\midrule
Phương sai ($\sigma^2$) & $X \sim N(\mu,\sigma^2)$, $\mu$ chưa biết & -- &
  $\dfrac{(n-1)s^2}{\chi^2_{\alpha/2;\,n-1}} < \sigma^2 < \dfrac{(n-1)s^2}{\chi^2_{1-\alpha/2;\,n-1}}$ \\
\bottomrule
\end{tabular}
\end{table}

\subsubsection{Kiểm định phân phối chuẩn}

Phân phối chuẩn là giả định cốt lõi của các kiểm định tham số (ANOVA, hồi quy OLS). Với
ANOVA, giả định này áp dụng cho biến phụ thuộc \emph{trong từng nhóm} (tương đương với phần
dư), không phải cho biến phụ thuộc gộp cả mẫu.
\begin{itemize}
  \item \textbf{Đồ thị xác suất chuẩn (Normal Q-Q Plot):} So sánh phân vị thực nghiệm của mẫu
  với phân vị lý thuyết của phân phối chuẩn tắc $N(0,1)$. Nếu các điểm bám sát đường thẳng
  45\textdegree, dữ liệu phù hợp với giả định phân phối chuẩn.
  \item \textbf{Kiểm định Shapiro-Wilk:} Cặp giả thuyết:
  \begin{equation}
    \begin{cases}
      H_0: \text{Dữ liệu tuân theo phân phối chuẩn}\\
      H_1: \text{Dữ liệu không tuân theo phân phối chuẩn}
    \end{cases}
  \end{equation}
  Thống kê kiểm định:
  \begin{equation}
    W = \frac{\left(\sum_{i=1}^{n} a_i x_{(i)}\right)^2}{\sum_{i=1}^{n}(x_i - \xbar)^2}
  \end{equation}
  trong đó $x_{(i)}$ là quan trắc thứ $i$ sau khi sắp xếp; $a_i$ là hệ số trọng số suy ra từ
  kỳ vọng và ma trận hiệp phương sai của các thống kê thứ tự của mẫu chuẩn kích thước $n$.
  Với $\alpha = 0{,}05$: nếu $p\text{-value} \ge 0{,}05$, chưa có cơ sở bác bỏ $H_0$; nếu
  $p\text{-value} < 0{,}05$, bác bỏ $H_0$. Lưu ý: với nhóm nhỏ (ví dụ $n = 6$), kiểm định có
  lực rất thấp, nên ``không bác bỏ $H_0$'' không chứng minh được dữ liệu chuẩn.
\end{itemize}

\subsection{Kiểm định so sánh trung bình ($t$-Test)}

\subsubsection{Kiểm định trung bình một mẫu (One-Sample $t$-Test)}

Áp dụng khi kiểm định giá trị trung bình tổng thể $\mu$ có bằng giá trị chuẩn $\mu_0$ khi chưa
biết phương sai tổng thể $\sigma^2$:
\begin{equation}
  T = \frac{\xbar - \mu_0}{s/\sqrt{n}} \sim t(n-1)
\end{equation}

\subsubsection{Kiểm định so sánh hai trung bình độc lập}

So sánh kỳ vọng hai tổng thể độc lập $\mu_1$ và $\mu_2$ từ hai mẫu kích thước $n_1$, $n_2$.
\begin{itemize}
  \item \textbf{Trường hợp phương sai bằng nhau ($\sigma_1^2 = \sigma_2^2$ -- Student's $t$-test):}
  \begin{equation}
    T = \frac{(\xbar_1 - \xbar_2) - D_0}{s_p\sqrt{\frac{1}{n_1}+\frac{1}{n_2}}} \sim t(n_1+n_2-2),
    \qquad s_p^2 = \frac{(n_1-1)s_1^2 + (n_2-1)s_2^2}{n_1+n_2-2}
  \end{equation}
  \item \textbf{Trường hợp phương sai khác nhau ($\sigma_1^2 \ne \sigma_2^2$ -- Welch's $t$-test):}
  \begin{equation}
    T = \frac{(\xbar_1 - \xbar_2) - D_0}{\sqrt{\frac{s_1^2}{n_1}+\frac{s_2^2}{n_2}}} \sim t(\nu)
  \end{equation}
  Bậc tự do hiệu chỉnh Welch--Satterthwaite:
  \begin{equation}
    \nu = \frac{\left(\frac{s_1^2}{n_1}+\frac{s_2^2}{n_2}\right)^2}
               {\frac{(s_1^2/n_1)^2}{n_1-1} + \frac{(s_2^2/n_2)^2}{n_2-1}}
  \end{equation}
\end{itemize}

\subsection{Phân tích phương sai một yếu tố (One-way ANOVA) \& Kiểm định phi tham số}

Phân tích phương sai một yếu tố dùng để kiểm định sự bằng nhau về giá trị trung bình của $k$
tổng thể độc lập ($k \ge 3$).

\subsubsection{Mô hình toán học}

Giả sử có $k$ nhóm độc lập, nhóm $i$ có $n_i$ quan trắc. Giá trị quan sát thứ $j$ thuộc nhóm
$i$ được mô hình hóa:
\begin{equation}
  Y_{ij} = \mu + \tau_i + \varepsilon_{ij} = \mu_i + \varepsilon_{ij}
  \qquad (i = 1,\dots,k;\; j = 1,\dots,n_i)
\end{equation}
trong đó $\mu$ là trung bình chung; $\mu_i$ là trung bình của nhóm $i$; $\tau_i = \mu_i - \mu$
là hiệu ứng của nhóm $i$; $\varepsilon_{ij}$ là sai số ngẫu nhiên.

\subsubsection{Các giả định bắt buộc của ANOVA}
\begin{enumerate}
  \item \textbf{Tính độc lập:} Các quan trắc trong từng nhóm và giữa các nhóm độc lập với nhau.
  \item \textbf{Tính phân phối chuẩn:} Biến phụ thuộc trong mỗi nhóm tuân theo phân phối chuẩn:
  $Y_{ij} \sim N(\mu_i, \sigma^2)$.
  \item \textbf{Tính đồng nhất phương sai:} $\sigma_1^2 = \sigma_2^2 = \dots = \sigma_k^2 = \sigma^2$.
\end{enumerate}

\subsubsection{Kiểm định tính đồng nhất phương sai (Levene -- biến thể Brown-Forsythe)}

Cặp giả thuyết: $H_0: \sigma_1^2 = \sigma_2^2 = \dots = \sigma_k^2$ đối
$H_1: \exists\,(i,l)$ sao cho $\sigma_i^2 \ne \sigma_l^2$. Thống kê được tính trên độ lệch tuyệt
đối so với \emph{trung vị} nhóm $Z_{ij} = |Y_{ij} - \tilde{Y}_i|$ (biến thể Brown-Forsythe), nên ít bị ảnh hưởng bởi tính phi chuẩn của dữ liệu:
\begin{equation}
  W = \frac{N-k}{k-1}\cdot
      \frac{\sum_{i=1}^{k} n_i(\bar{Z}_{i\cdot} - \bar{Z}_{\cdot\cdot})^2}
           {\sum_{i=1}^{k}\sum_{j=1}^{n_i}(Z_{ij} - \bar{Z}_{i\cdot})^2}
  \sim F(k-1,\, N-k)
\end{equation}

\subsubsection{Phân rã tổng các bình phương ($SST = SSB + SSW$) \& Bảng ANOVA}

Tổng biến thiên toàn phần ($SST$) được phân rã thành biến thiên giữa các nhóm ($SSB$) và biến
thiên nội bộ nhóm ($SSW$):
\begin{equation}
  SST = SSB + SSW
\end{equation}
\begin{itemize}
  \item Tổng bình phương toàn phần:
  \begin{equation}
    SST = \sum_{i=1}^{k}\sum_{j=1}^{n_i}(Y_{ij} - \bar{Y}_{\cdot\cdot})^2, \quad df_T = N-1
  \end{equation}
  \item Tổng bình phương giữa các nhóm:
  \begin{equation}
    SSB = \sum_{i=1}^{k} n_i(\bar{Y}_{i\cdot} - \bar{Y}_{\cdot\cdot})^2, \quad df_B = k-1
  \end{equation}
  \item Tổng bình phương trong nội bộ nhóm:
  \begin{equation}
    SSW = \sum_{i=1}^{k}\sum_{j=1}^{n_i}(Y_{ij} - \bar{Y}_{i\cdot})^2, \quad df_W = N-k
  \end{equation}
\end{itemize}
(Trong đó $N = \sum_{i=1}^{k} n_i$; $\bar{Y}_{i\cdot}$ là trung bình nhóm $i$;
$\bar{Y}_{\cdot\cdot}$ là trung bình chung toàn bộ mẫu.)

\begin{table}[H]
\centering
\caption{Bảng phân tích phương sai chuẩn (ANOVA)}
\label{tab:anova}
\begin{tabular}{l c c c c c}
\toprule
Nguồn biến thiên & SS & df & MS & Thống kê $F$ & $p$-value \\
\midrule
Giữa các nhóm & $SSB$ & $k-1$ & $MSB = \frac{SSB}{k-1}$ & $F = \frac{MSB}{MSW}$ & $P(F_{k-1,N-k} > F)$ \\
Nội bộ nhóm   & $SSW$ & $N-k$ & $MSW = \frac{SSW}{N-k}$ & & \\
Toàn phần     & $SST$ & $N-1$ & & & \\
\bottomrule
\end{tabular}
\end{table}

Cặp giả thuyết kiểm định:
\begin{equation}
  \begin{cases}
    H_0: \mu_1 = \mu_2 = \dots = \mu_k\\
    H_1: \exists\,(i,l) \text{ sao cho } \mu_i \ne \mu_l
  \end{cases}
\end{equation}
Quy tắc kết luận: Nếu $F > F_{\alpha;\,k-1,\,N-k}$ (hoặc $p\text{-value} < \alpha$), bác bỏ $H_0$
và kết luận có ít nhất một cặp nhóm có trung bình khác biệt. Một nhóm chỉ có 1 quan trắc vẫn
tính được $F$ (góp 0 bậc tự do vào $SSW$), nhưng không kiểm tra được giả định cho nhóm đó.

\subsubsection{Phân tích so sánh bội hậu nghiệm (Tukey HSD)}

Khi ANOVA bác bỏ $H_0$, kiểm định Tukey HSD được dùng để kiểm soát sai lầm loại I toàn họ phép
thử (familywise error rate) khi so sánh từng cặp trung bình $(\mu_i, \mu_l)$. Khoảng tin cậy đồng
thời cho hiệu $(\mu_i - \mu_l)$:
\begin{equation}
  (\bar{Y}_{i\cdot} - \bar{Y}_{l\cdot}) \pm q_{\alpha;\,k,\,N-k}
  \sqrt{\frac{MSW}{2}\left(\frac{1}{n_i}+\frac{1}{n_l}\right)}
\end{equation}
Nếu khoảng tin cậy không chứa giá trị 0, sự khác biệt giữa hai nhóm có ý nghĩa thống kê.

\subsubsection{Kiểm định phi tham số thay thế: Kruskal-Wallis Test}

Áp dụng khi dữ liệu vi phạm giả định chuẩn. Xếp hạng toàn bộ $N$ quan sát (hạng trùng nhận hạng
trung bình), gọi $R_i$ là tổng hạng của nhóm $i$:
\begin{equation}
  H = \frac{12}{N(N+1)}\sum_{i=1}^{k}\frac{R_i^2}{n_i} - 3(N+1)
\end{equation}
Khi có hạng trùng, thống kê được hiệu chỉnh:
\begin{equation}
  H_{c} = \frac{H}{1 - \dfrac{\sum_{g}(t_g^3 - t_g)}{N^3 - N}} \;\approx\; \chi^2(k-1)
\end{equation}
với $t_g$ là số quan sát trong nhóm hạng trùng thứ $g$. Khi $H_c$ có ý nghĩa thống kê, kiểm
định hậu nghiệm Dunn (kèm hiệu chỉnh Holm hoặc Bonferroni) được dùng để so sánh từng cặp nhóm.

\subsection{Phân tích tương quan}

\subsubsection{Hệ số tương quan Pearson ($r$) \& Kiểm định ý nghĩa}

Đo lường mức độ và chiều hướng tương quan tuyến tính giữa hai biến định lượng $X$ và $Y$:
\begin{equation}
  r = \frac{\sum_{i=1}^{n}(x_i - \xbar)(y_i - \bar{y})}
           {\sqrt{\sum_{i=1}^{n}(x_i - \xbar)^2}\sqrt{\sum_{i=1}^{n}(y_i - \bar{y})^2}}
\end{equation}
Tính chất: $-1 \le r \le 1$. Quy ước: $|r| \ge 0{,}7$: tương quan mạnh;
$0{,}4 \le |r| < 0{,}7$: tương quan trung bình; $|r| < 0{,}4$: tương quan yếu.
Hệ số Pearson nhạy với giá trị ngoại lai: một vài điểm ở xa có thể làm $|r|$ tăng hoặc giảm
đáng kể.

Kiểm định ý nghĩa thống kê của hệ số tương quan ($H_0: \rho = 0$ đối $H_1: \rho \ne 0$):
\begin{equation}
  T = \frac{r\sqrt{n-2}}{\sqrt{1-r^2}} \sim t(n-2)
\end{equation}

\subsubsection{Hệ số tương quan hạng Spearman ($r_s$)}

Đo lường mối quan hệ đơn điệu (không nhất thiết tuyến tính), không yêu cầu phân phối chuẩn và
vững hơn trước ngoại lai. Về định nghĩa, $r_s$ là hệ số Pearson tính trên thứ hạng của hai
biến (hạng trùng nhận hạng trung bình). Khi \emph{không có hạng trùng},
công thức trên rút gọn thành:
\begin{equation}
  r_s = 1 - \frac{6\sum_{i=1}^{n} d_i^2}{n(n^2-1)}
\end{equation}
với $d_i = \operatorname{rank}(x_i) - \operatorname{rank}(y_i)$. Dữ liệu của nhóm có nhiều giá
trị trùng (calories, protein, vitamins\ldots), nên công thức rút gọn chỉ mang tính minh họa.

\subsection{Mô hình Hồi quy tuyến tính bội (Multiple Linear Regression -- OLS)}
\label{sec:ols}

\subsubsection{Dạng tổng quát mô hình}
\begin{itemize}
  \item Dạng đại số:
  \begin{equation}
    Y_i = \beta_0 + \beta_1 X_{i1} + \beta_2 X_{i2} + \dots + \beta_p X_{ip} + \varepsilon_i
    \qquad (i = 1,\dots,n)
  \end{equation}
  \item Dạng ma trận: $\mathbf{Y} = \mathbf{X}\boldsymbol{\beta} + \boldsymbol{\varepsilon}$, với
  \begin{equation}
    \mathbf{Y} = \begin{pmatrix} Y_1\\ Y_2\\ \vdots\\ Y_n \end{pmatrix},\;
    \mathbf{X} = \begin{pmatrix}
      1 & X_{11} & \cdots & X_{1p}\\
      1 & X_{21} & \cdots & X_{2p}\\
      \vdots & \vdots & \ddots & \vdots\\
      1 & X_{n1} & \cdots & X_{np}
    \end{pmatrix},\;
    \boldsymbol{\beta} = \begin{pmatrix} \beta_0\\ \beta_1\\ \vdots\\ \beta_p \end{pmatrix},\;
    \boldsymbol{\varepsilon} = \begin{pmatrix} \varepsilon_1\\ \varepsilon_2\\ \vdots\\ \varepsilon_n \end{pmatrix}
  \end{equation}
\end{itemize}

\subsubsection{Phương pháp bình phương bé nhất (OLS)}

Ước lượng $\hat{\boldsymbol{\beta}}$ được tìm bằng cách cực tiểu hóa tổng bình phương phần dư
($SSE$):
\begin{equation}
  \hat{\boldsymbol{\beta}} = (\mathbf{X}^\top\mathbf{X})^{-1}\mathbf{X}^\top\mathbf{Y}
\end{equation}
Giá trị dự báo: $\hat{\mathbf{Y}} = \mathbf{X}\hat{\boldsymbol{\beta}} = \mathbf{H}\mathbf{Y}$
với $\mathbf{H} = \mathbf{X}(\mathbf{X}^\top\mathbf{X})^{-1}\mathbf{X}^\top$ là ma trận Hat.
Vector phần dư: $\mathbf{e} = \mathbf{Y} - \hat{\mathbf{Y}} = (\mathbf{I} - \mathbf{H})\mathbf{Y}$.

\subsubsection{Đánh giá mức độ phù hợp của mô hình}

Phân rã biến thiên toàn phần:
\begin{equation}
  SST = SSR + SSE \iff
  \sum_{i=1}^{n}(Y_i - \bar{Y})^2 = \sum_{i=1}^{n}(\hat{Y}_i - \bar{Y})^2 + \sum_{i=1}^{n}(Y_i - \hat{Y}_i)^2
\end{equation}
\begin{itemize}
  \item Hệ số xác định:
  \begin{equation}
    R^2 = \frac{SSR}{SST} = 1 - \frac{SSE}{SST}
  \end{equation}
  \item Hệ số xác định hiệu chỉnh:
  \begin{equation}
    R^2_{\text{adj}} = 1 - \frac{SSE/(n-p-1)}{SST/(n-1)} = 1 - (1-R^2)\frac{n-1}{n-p-1}
  \end{equation}
\end{itemize}

\subsubsection{Các kiểm định ý nghĩa trong mô hình hồi quy}
\begin{itemize}
  \item \textbf{Kiểm định độ phù hợp chung của mô hình (Kiểm định $F$):}
  $H_0: \beta_1 = \beta_2 = \dots = \beta_p = 0$ đối $H_1: \exists\,\beta_j \ne 0$.
  \begin{equation}
    F = \frac{SSR/p}{SSE/(n-p-1)} \sim F(p,\, n-p-1)
  \end{equation}
  Nếu $p\text{-value} < \alpha$, mô hình có ý nghĩa giải thích tổng thể.
  \item \textbf{Kiểm định ý nghĩa của từng hệ số hồi quy (Kiểm định $t$):}
  $H_0: \beta_j = 0$ đối $H_1: \beta_j \ne 0$.
  \begin{equation}
    T_j = \frac{\hat{\beta}_j}{SE(\hat{\beta}_j)} \sim t(n-p-1)
  \end{equation}
  với $SE(\hat{\beta}_j) = \sqrt{\hat{\sigma}^2\,[(\mathbf{X}^\top\mathbf{X})^{-1}]_{jj}}$ và
  $\hat{\sigma}^2 = \dfrac{SSE}{n-p-1}$.
\end{itemize}

\subsubsection{Kiểm định và chẩn đoán các giả định của sai số OLS}
\begin{enumerate}
  \item \textbf{Đa cộng tuyến (Multicollinearity):} Đánh giá qua hệ số phóng đại phương sai:
  \begin{equation}
    VIF_j = \frac{1}{1 - R_j^2}
  \end{equation}
  trong đó $R_j^2$ là hệ số xác định khi hồi quy biến độc lập $X_j$ theo toàn bộ các biến độc
  lập còn lại. Ngưỡng $VIF < 5$ là an toàn; $VIF \ge 5$ cần chú ý và $VIF \ge 10$ cảnh báo đa
  cộng tuyến nghiêm trọng cần khắc phục.
  \item \textbf{Tính tuyến tính và kỳ vọng sai số bằng 0 ($E(\varepsilon) = 0$):} Kiểm tra qua đồ
  thị phần dư theo giá trị dự báo (Residuals vs Fitted). Các điểm phải phân tán ngẫu nhiên xung
  quanh đường $e = 0$.
  \item \textbf{Phương sai sai số đồng nhất ($\operatorname{Var}(\varepsilon) = \sigma^2$):} Kiểm tra
  qua biểu đồ Scale-Location và kiểm định Breusch-Pagan với $H_0$: phương sai sai số không đổi.
  \item \textbf{Sai số tuân theo phân phối chuẩn ($\varepsilon \sim N(0,\sigma^2)$):} Kiểm tra qua
  đồ thị Normal Q-Q của phần dư chuẩn hóa và kiểm định Shapiro-Wilk trên vector phần dư $\mathbf{e}$.
  \item \textbf{Tính độc lập của sai số:} Kiểm định Durbin-Watson:
  \begin{equation}
    DW = \frac{\sum_{i=2}^{n}(e_i - e_{i-1})^2}{\sum_{i=1}^{n} e_i^2}
  \end{equation}
  $DW \approx 2$ cho thấy không có tự tương quan \emph{bậc 1} giữa các sai số liền kề; kiểm
  định này không kiểm tra các dạng phụ thuộc khác. Với dữ liệu chéo như bộ dữ liệu này, thứ tự
  dòng là thứ tự tên sản phẩm theo bảng chữ cái, nên kết quả DW gần như không mang ý nghĩa và
  chỉ được báo cáo để tham khảo.
  \item \textbf{Nhận diện điểm ảnh hưởng cực đoan:} Sử dụng khoảng cách Cook:
  \begin{equation}
    D_i = \frac{\sum_{j=1}^{n}(\hat{Y}_j - \hat{Y}_{j(i)})^2}{(p+1)\,\hat{\sigma}^2}
  \end{equation}
  Các quan trắc có $D_i > 0{,}5$ hoặc $D_i > 4/n$ được xem là điểm có ảnh hưởng bất thường cần
  xem xét.
\end{enumerate}

% =====================================================================
\section{TIỀN XỬ LÝ SỐ LIỆU}
% =====================================================================

Quy trình tiền xử lý được thực hiện trên phần mềm R theo đúng thứ tự các mục 4.1 -- 4.5 dưới đây.

\subsection{Đọc dữ liệu và Khảo sát dữ liệu khuyết}

Tập dữ liệu \texttt{cereal.csv} gồm 77 quan sát và 16 biến. Dữ liệu không có ô rỗng
\texttt{NA} nhưng xuất hiện mã ẩn $-1$ đại diện cho giá trị khuyết tại 3 biến định lượng:
\texttt{potass} (2 dòng: 2,6\%, gồm Almond Delight và Cream of Wheat (Quick)), \texttt{carbo}
(1 dòng: 1,3\%) và \texttt{sugars} (1 dòng: 1,3\%), hai giá trị sau cùng thuộc Quaker
Oatmeal. Nhóm chuyển đổi các giá trị $-1$ này thành \texttt{NA} chuẩn trong R để chuẩn bị cho
bước xử lý tiếp theo.

\begin{lstlisting}
# 4.1. Đọc dữ liệu gốc và chuyển mã ẩn -1 thành NA
num_cols <- c("calories", "protein", "fat", "sodium", "fiber",
              "carbo", "sugars", "potass", "weight", "cups", "rating")

read_raw_cereal <- function(path = "cereal.csv") {
  cereal <- read.csv(path, stringsAsFactors = FALSE)
  cereal[num_cols] <- lapply(cereal[num_cols], function(x) {
    ifelse(x == -1, NA, as.numeric(x))
  })
  cereal
}
\end{lstlisting}

\subsection{Lọc mẫu và Bớt biến sơ bộ}

Việc lọc mẫu được thực hiện \emph{trước} bước điền khuyết, để trung vị dùng điền khuyết được
tính trên đúng mẫu phân tích $N = 76$. Cụ thể, nhóm loại bỏ 1 quan trắc thuộc nhà sản xuất A (American Home Food
Products, sản phẩm Maypo) vì $n_A = 1$ không cho phép kiểm tra giả định và so sánh cặp cho
nhóm này, giữ lại tập mẫu phân tích chung gồm $N = 76$ quan trắc.

Song song đó, nhóm loại bỏ biến \texttt{name} (chỉ có chức năng định danh) và
biến \texttt{type} (chỉ có 3/77 sản phẩm ăn nóng). Tên sản phẩm được giữ làm tên dòng để tra
cứu ngoại lai ở Phần~\ref{sec:mota}. Biến \texttt{weight} được giữ lại tạm thời để dùng cho
bước quy đổi khẩu phần.

\begin{lstlisting}
# 4.2. Lọc bỏ hãng A (n = 1) trước khi điền khuyết; bỏ name, type.
# Tên sản phẩm được giữ làm tên dòng để tra cứu ngoại lai ở Phần 5.
filter_sample <- function(cereal) {
  out <- subset(cereal, mfr != "A")
  rownames(out) <- out$name
  out[, !(names(out) %in% c("name", "type"))]
}
\end{lstlisting}

\subsection{Điền khuyết dữ liệu}

Trên tập dữ liệu $N = 76$, nhóm điền các giá trị khuyết bằng trung vị của từng biến tương ứng.
Phương pháp điền trung vị là thước đo vững, ít làm méo phân bố khi tỷ lệ khuyết rất nhỏ
($< 3\%$) và bảo toàn cỡ mẫu $N = 76$. Hạn chế của cách làm này là giá trị điền có thể không
thực tế với sản phẩm cụ thể: Quaker Oatmeal là yến mạch nóng không đường nhưng được điền
\texttt{sugars} = 7 g và \texttt{carbo} = 14 g. Do chỉ ảnh hưởng 3 dòng, nhóm chấp nhận hạn chế
này cho phần mô tả và ghi nhận lại để xem xét ở phần suy diễn.

\begin{lstlisting}
# 4.3. Điền giá trị khuyết (NA) bằng trung vị tính trên mẫu N = 76
impute_median <- function(cereal, cols = c("carbo", "sugars", "potass")) {
  for (col in cols) {
    med_val <- median(cereal[[col]], na.rm = TRUE)
    cereal[[col]][is.na(cereal[[col]])] <- med_val
  }
  cereal
}
\end{lstlisting}

\subsection{Chuẩn hóa quy mô khẩu phần và Tạo biến mới}

Do khối lượng khẩu phần (\texttt{weight}) giữa các loại ngũ cốc không đồng nhất (từ 0,5 oz
đến 1,5 oz), so sánh trực tiếp các thành phần dinh dưỡng nguyên bản sẽ gây sai lệch do quy mô
khẩu phần khác nhau. Vì vậy, nhóm quy đổi 8 biến dinh dưỡng về hàm lượng trên 1 oz
bằng cách chia giá trị gốc cho \texttt{weight}:
\begin{equation}
\text{Nutrient}_{\text{norm}} = \frac{\text{Nutrient}}{\text{weight}}
\end{equation}
\texttt{rating} và \texttt{vitamins} giữ nguyên. Lưu ý theo mục~\ref{sec:rating}:
\texttt{rating} được tính theo khẩu phần gốc, nên với 13 sản phẩm có weight $\neq 1$, giá
trị dinh dưỡng sau chuẩn hóa không còn khớp chính xác với cách rating được tính.

Sau khi quy đổi, nhóm loại bỏ \texttt{weight} và
\texttt{cups} khỏi tập dữ liệu phân tích.

Cuối cùng, nhóm tạo biến khám phá \texttt{total\_carbs} = \texttt{carbo} +
\texttt{sugars} (tổng lượng bột đường trên 1 oz), chỉ phục vụ Thống kê mô tả.

\begin{lstlisting}
# 4.4. Chuẩn hóa 8 biến dinh dưỡng về khẩu phần 1 oz, tạo total_carbs,
# rồi bỏ weight và cups
nutri_cols <- c("calories", "protein", "fat", "sodium",
                "fiber", "carbo", "sugars", "potass")

normalize_per_oz <- function(cereal, cols = nutri_cols) {
  for (col in cols) {
    cereal[[col]] <- cereal[[col]] / cereal$weight
  }
  cereal$total_carbs <- cereal$carbo + cereal$sugars
  cereal[, !(names(cereal) %in% c("weight", "cups"))]
}
\end{lstlisting}

\subsection{Định dạng biến Factor và Kiểm tra cấu trúc dữ liệu}

Nhóm chuyển các biến phân nhóm (\texttt{mfr}, \texttt{shelf}) từ dạng ký tự/số nguyên sang
dạng Factor có nhãn rõ ràng để sẵn sàng cho các kiểm định suy diễn (ANOVA, Kruskal-Wallis) và
hồi quy. Biến \texttt{vitamins} được giữ nguyên dạng số. Dữ liệu sạch (76 dòng $\times$ 13 cột)
được lưu lại để sử dụng cho các phần phân tích tiếp theo. Nhóm kiểm tra lại dữ liệu sạch theo 10
điều kiện (số dòng, số cột, không còn \texttt{NA}, đã chia \texttt{weight}, điền đúng trung
vị\ldots) và đều đạt.

\begin{lstlisting}
# 4.5. Đổi mfr và shelf sang factor có nhãn rõ ràng (vitamins giữ dạng số)
mfr_labels <- c("G" = "General Mills", "K" = "Kelloggs", "N" = "Nabisco",
                "P" = "Post", "Q" = "Quaker Oats", "R" = "Ralston Purina")
shelf_labels <- c("Kệ 1 (Thấp)", "Kệ 2 (Vừa)", "Kệ 3 (Cao)")

format_factors <- function(cereal) {
  cereal$mfr <- factor(mfr_labels[cereal$mfr])
  cereal$shelf <- factor(cereal$shelf, levels = c(1, 2, 3), labels = shelf_labels)
  cereal
}

# Toàn bộ Phần 4: dữ liệu gốc (77 x 16) -> dữ liệu sạch (76 x 13)
cereal_raw   <- read_raw_cereal()
cereal_final <- format_factors(normalize_per_oz(impute_median(filter_sample(cereal_raw))))
write.csv(cereal_final, "cereal_cleaned.csv", row.names = FALSE)
\end{lstlisting}

% =====================================================================
\section{THỐNG KÊ MÔ TẢ}
\label{sec:mota}
% =====================================================================

\subsection{Tính toán các thống kê mô tả}

\subsubsection{Thống kê mô tả biến định tính (mfr, shelf) và biến rời rạc (vitamins)}

\begin{table}[H]
\centering
\caption{Phân bố tần số, tỷ lệ phần trăm của nhà sản xuất, vị trí kệ và hàm lượng vitamin}
\label{tab:tanso}
\begin{tabular}{l l c c}
\toprule
Biến & Phân nhóm & Tần số ($n_k$) & Tỷ lệ ($f_k$, \%) \\
\midrule
\multirow{6}{*}{Nhà sản xuất (mfr)}
  & General Mills  & 22 & 28,95 \\
  & Kelloggs       & 23 & 30,26 \\
  & Nabisco        & 6  & 7,89 \\
  & Post           & 9  & 11,84 \\
  & Quaker Oats    & 8  & 10,53 \\
  & Ralston Purina & 8  & 10,53 \\
\midrule
\multirow{3}{*}{Vị trí kệ (shelf)}
  & Kệ 1 (Thấp) & 20 & 26,32 \\
  & Kệ 2 (Vừa)  & 20 & 26,32 \\
  & Kệ 3 (Cao)  & 36 & 47,37 \\
\midrule
\multirow{3}{*}{Hàm lượng vitamin (\%)}
  & 0\%   & 8  & 10,53 \\
  & 25\%  & 62 & 81,58 \\
  & 100\% & 6  & 7,89 \\
\bottomrule
\end{tabular}
\end{table}

Từ kết quả Bảng~\ref{tab:tanso}, cơ cấu mẫu bộc lộ ba đặc điểm phân bố đáng chú ý.

\textbf{Sự mất cân bằng cỡ mẫu của biến mfr:} Mẫu bị chi phối bởi Kelloggs và General
Mills, gộp chung gần 60\% tổng số sản phẩm (59,21\%). Ngược lại, các hãng còn lại có quy mô
mẫu khá nhỏ, đặc biệt là Nabisco chỉ có 6 sản phẩm. Sự chênh lệch cỡ mẫu này đòi hỏi phải
kiểm tra chặt chẽ tính đồng nhất phương sai, và kết quả ở các nhóm nhỏ cần được diễn giải
thận trọng.

\textbf{Sự tập trung trưng bày tại vị trí kệ cao:} Kệ 3 (Cao) là nhóm áp đảo với 36 sản
phẩm (chiếm 47,37\%), gần gấp đôi quy mô của Kệ 1 (Thấp) và Kệ 2 (Vừa).

\textbf{Sự co cụm mạnh của vitamins:} Có tới 62/76 sản phẩm được bổ sung vi chất ở mức
tiêu chuẩn (25\%) nhu cầu khuyến nghị hằng ngày, trong khi nhóm không bổ sung (0\%) và nhóm
tăng cường tối đa (100\%) chỉ chiếm tỷ trọng nhỏ.

\subsubsection{Thống kê mô tả các biến định lượng}

\begin{table}[H]
\centering
\caption{Thống kê mô tả các biến định lượng (8 biến dinh dưỡng và total\_carbs tính trên 1 oz; rating và vitamins giữ nguyên)}
\label{tab:dinhluong}
\small
\setlength{\tabcolsep}{4pt}
\begin{tabular}{l c r r r r r r r r r}
\toprule
Biến & $n$ & Mean & SD & Min & $Q_1$ & Median & $Q_3$ & Max & IQR & Skewness \\
\midrule
rating      & 76 & 42,51  & 14,07 & 18,04 & 32,93  & 40,25  & 50,78  & 93,70  & 17,85 & 0,94 \\
calories    & 76 & 104,05 & 13,94 & 50,00 & 100,00 & 107,18 & 110,00 & 150,00 & 10,00 & $-0{,}22$ \\
protein     & 76 & 2,47   & 1,08  & 1,00  & 2,00   & 2,13   & 3,00   & 6,00   & 1,00  & 1,03 \\
fat         & 76 & 0,97   & 0,99  & 0,00  & 0,00   & 1,00   & 1,38   & 5,00   & 1,38  & 1,33 \\
sodium      & 76 & 155,37 & 81,39 & 0,00  & 126,25 & 161,65 & 202,50 & 320,00 & 76,25 & $-0{,}41$ \\
fiber       & 76 & 2,08   & 2,31  & 0,00  & 1,00   & 1,77   & 3,00   & 14,00  & 2,00  & 2,72 \\
carbo       & 76 & 14,60  & 4,23  & 5,00  & 11,46  & 14,00  & 17,00  & 26,00  & 5,54  & 0,35 \\
sugars      & 76 & 6,71   & 4,07  & 0,00  & 3,00   & 6,47   & 10,00  & 15,00  & 7,00  & 0,08 \\
potass      & 76 & 93,75  & 63,22 & 20,00 & 43,75  & 90,00  & 114,59 & 330,00 & 70,84 & 1,66 \\
total\_carbs & 76 & 21,31 & 3,48  & 8,00  & 19,55  & 21,00  & 24,00  & 27,00  & 4,45  & $-1{,}13$ \\
vitamins    & 76 & 28,29  & 22,49 & 0,00  & 25,00  & 25,00  & 25,00  & 100,00 & 0,00  & 2,44 \\
\bottomrule
\end{tabular}
\end{table}

Căn cứ vào hệ số bất đối xứng $G_1$ theo quy ước ở mục 3.1.2 và sự chênh lệch giữa trung bình
($\xbar$) với trung vị ($\xmed$), 11 biến định lượng được phân thành 3 nhóm hình dạng phân phối.

\textbf{Nhóm lệch mạnh ($|G_1| > 1$):} Lệch phải gồm fiber (2,72), vitamins (2,44),
potass (1,66), fat (1,33) và protein (1,03); phần lớn sản phẩm tập trung ở vùng giá trị thấp,
trong khi một số ít ngũ cốc chuyên biệt kéo dài đuôi phân phối về phía bên phải. Riêng
\texttt{total\_carbs} lệch trái mạnh ($G_1 = -1{,}13$): phần lõi phân phối ổn định
($IQR = 4{,}45$ g quanh trung vị $\xmed = 21{,}00$ g) nhưng có đuôi trái do các dòng ngũ cốc
cám ít tinh bột.

\textbf{Nhóm lệch vừa ($0{,}5 \le |G_1| \le 1$):} Biến phụ thuộc \texttt{rating} lệch
phải vừa ($G_1 = 0{,}94$; $\xbar = 42{,}51 > \xmed = 40{,}25$ điểm; khoảng biến thiên rộng từ
18,04 đến 93,70 điểm).

\textbf{Nhóm tương đối đối xứng ($|G_1| < 0{,}5$):} Gồm sugars (0,08), carbo (0,35),
calories ($-0{,}22$) và sodium ($-0{,}41$). Hàm lượng đường và tinh bột phân bố khá cân đối
quanh tâm. Calories có hệ số bất đối xứng nhỏ nhưng co cụm rất chặt quanh vùng
100 -- 110 kcal/oz ($IQR = 10{,}00$ kcal) và có ngoại lai ở cả hai phía; sodium hơi lệch trái
nhẹ do nhóm sản phẩm không thêm muối (0 mg).

\subsubsection{Phát hiện và biện luận giá trị ngoại lai theo hàng rào Tukey}

\begin{table}[H]
\centering
\caption{Tổng hợp giá trị ngoại lai theo hàng rào Tukey và sản phẩm tương ứng}
\label{tab:ngoailai}
\small
\begin{tabular}{l r r c L{7.6cm}}
\toprule
Biến & \makecell[r]{Ngưỡng\\dưới ($L$)} & \makecell[r]{Ngưỡng\\trên ($U$)} & \makecell{Số\\lượng} & Giá trị ngoại lai và tên sản phẩm cụ thể \\
\midrule
rating & 6,16 & 77,55 & 1 & 93,70 (All-Bran with Extra Fiber) \\
\midrule
calories & 85,00 & 125,00 & 5 & 50,00 (All-Bran with Extra Fiber); 70,00 (100\% Bran);
  70,00 (All-Bran); 150,00 (Muesli Raisins, Dates, \& Almonds); 150,00 (Muesli Raisins,
  Peaches, \& Pecans) \\
\midrule
protein & 0,50 & 4,50 & 3 & 5,00 (Quaker Oatmeal); 6,00 (Cheerios); 6,00 (Special K) \\
\midrule
fat & $-2{,}06$ & 3,44 & 1 & 5,00 (100\% Natural Bran) \\
\midrule
sodium & 11,88 & 316,88 & 9 & 0,00 (Frosted Mini-Wheats, Puffed Rice, Puffed Wheat, Quaker
  Oatmeal, Raisin Squares, Shredded Wheat, Shredded Wheat 'n'Bran, Shredded Wheat spoon size);
  320,00 (Product 19) \\
\midrule
fiber & $-2{,}00$ & 6,00 & 3 & 9,00 (All-Bran); 10,00 (100\% Bran); 14,00 (All-Bran with Extra Fiber) \\
\midrule
carbo & 3,15 & 25,31 & 1 & 26,00 (Puffed Rice) \\
\midrule
potass & $-62{,}52$ & 220,86 & 3 & 280,00 (100\% Bran); 320,00 (All-Bran); 330,00 (All-Bran with Extra Fiber) \\
\midrule
total\_carbs & 12,87 & 30,68 & 3 & 8,00 (All-Bran with Extra Fiber); 11,00 (100\% Bran); 12,00 (All-Bran) \\
\midrule
vitamins & 25,00 & 25,00 & 14 & 0,00 (8 sản phẩm mức 0\%); 100,00 (6 sản phẩm mức 100\%) \\
\bottomrule
\end{tabular}
\end{table}

Từ Bảng~\ref{tab:ngoailai}, nhóm rút ra các nhận định và phương án xử lý ngoại lai như sau.

\textbf{Cụm ngũ cốc giàu cám lúa mì:} Các ngoại lai trên của fiber, potass, rating cùng
các ngoại lai dưới của calories và total\_carbs đều rơi đồng thời vào đúng 3 sản phẩm chuyên
biệt (All-Bran with Extra Fiber, All-Bran, 100\% Bran). Đây là đặc tính dinh dưỡng thực tế
của dòng sản phẩm giàu chất xơ -- ít năng lượng, không phải sai số đo lường.

\textbf{Các ngoại lai đặc thù sản phẩm khác:} 8 ngoại lai dưới ở sodium (0,00 mg) đại
diện cho các dòng ngũ cốc nguyên hạt không thêm muối (như Shredded Wheat, Puffed Rice, Quaker
Oatmeal), còn 14 điểm ngoại lai ở vitamins phát sinh do $IQR = Q_3 - Q_1 = 25 - 25 = 0$ khi
81,58\% mẫu tập trung tại mức 25\%; đây là hệ quả của biến rời rạc, không phải giá trị bất
thường.

\textbf{Phương án xử lý:} Toàn bộ các điểm ngoại lai đều phản ánh đúng bản chất đa
dạng của sản phẩm thương mại. Do đó, nhóm quyết định giữ nguyên toàn bộ $N = 76$ quan sát,
đồng thời ưu tiên các thước đo vững (Median, IQR, hệ số Spearman) và các kiểm định phi tham số
khi cần thiết. Ảnh hưởng của cụm 3 sản phẩm cám lên hệ số tương quan được đánh giá cụ thể ở
mục~\ref{sec:tuongquan}.

\subsection{Mô tả dữ liệu bằng đồ thị}

\subsubsection{Biểu đồ phân phối tần suất (Histogram kết hợp đường mật độ KDE)}

\begin{figure}[H]
\centering
\includegraphics[width=\textwidth]{Hinh_1_Histogram_KDE.png}
\caption{Phân phối tần suất (Histogram + KDE) của 10 biến định lượng liên tục (N = 76). Nét
đứt đỏ: trung bình; nét liền xanh: trung vị.}
\label{fig:hist}
\end{figure}

Đối chiếu trực quan trên Hình~\ref{fig:hist} cho thấy sự tương đồng với các chỉ số tại
Bảng~\ref{tab:dinhluong}. Các biến fiber, potass, protein và biến phụ thuộc rating có đỉnh phân phối lệch hẳn về
bên trái, đường trung bình (nét đứt đỏ) nằm bên phải đường trung vị (nét liền xanh) và kéo
đuôi dài về phía giá trị cao. Biến fat cũng lệch phải ($G_1 = 1{,}33$) nhưng là ngoại lệ về vị
trí hai đường tham chiếu: trung bình $\xbar = 0{,}97$ nhỏ hơn trung vị $\xmed = 1{,}00$, do
có nhiều sản phẩm 0 g chất béo kéo trung bình xuống.
Ngược lại, sugars và carbo có dạng phân phối trải đều và cân xứng hơn quanh hai đường
tham chiếu trung tâm, còn calories và total\_carbs tập trung dày đặc quanh vùng đỉnh
100 -- 110 kcal và 20 -- 24 g.
Hình dạng lệch của nhiều biến gợi ý rằng giả định phân phối chuẩn có thể không thỏa.
Tuy nhiên, với ANOVA, giả định chuẩn áp dụng cho rating \emph{trong từng nhóm} (hoặc phần dư),
nên kết luận chỉ được rút ra sau khi kiểm định Shapiro-Wilk theo từng nhóm ở phần suy diễn;
histogram của toàn mẫu chỉ mang tính gợi ý.

\subsubsection{Biểu đồ hộp đơn biến (Boxplot)}

\begin{figure}[H]
\centering
\includegraphics[width=\textwidth]{Hinh_2_Boxplot.png}
\caption{Biểu đồ hộp (Boxplot) phát hiện giá trị ngoại lai theo hàng rào Tukey (N = 76)}
\label{fig:box}
\end{figure}

Biểu đồ hộp nằm ngang ở Hình~\ref{fig:box} xác nhận trực quan vị trí và mật độ các điểm dị
biệt (tô đỏ) khớp với kết quả tính toán tại Bảng~\ref{tab:ngoailai}. Các biến fiber, potass, protein, fat và rating chỉ xuất hiện ngoại lai phía trên râu
phải (phản ánh nhóm sản phẩm giàu dưỡng chất), trong khi total\_carbs có ngoại lai ở đầu trái
và carbo có ngoại lai ở đầu phải.
Các biến calories, sodium và vitamins có ngoại lai ở cả hai phía; riêng sugars có hộp tứ
phân vị cân đối và hoàn toàn không có điểm nào vượt khỏi hai đầu râu Tukey.

\subsubsection{Biểu đồ hộp phân nhóm (Grouped Boxplot) theo mfr và shelf}

\begin{table}[H]
\centering
\caption{Thống kê mô tả điểm đánh giá (rating) phân tổ theo Nhà sản xuất và Vị trí kệ}
\label{tab:nhom}
\begin{tabular}{l l c r r r r}
\toprule
Yếu tố phân nhóm & Phân nhóm & $n$ & Mean & SD & Median & IQR \\
\midrule
\multirow{6}{*}{Nhà sản xuất (mfr)}
  & General Mills  & 22 & 34,49 & 8,95  & 36,18 & 11,62 \\
  & Kelloggs       & 23 & 44,04 & 14,46 & 40,56 & 15,54 \\
  & Nabisco        & 6  & 67,97 & 5,51  & 68,32 & 6,24 \\
  & Post           & 9  & 41,71 & 10,05 & 40,92 & 16,82 \\
  & Quaker Oats    & 8  & 42,92 & 16,80 & 47,42 & 22,35 \\
  & Ralston Purina & 8  & 41,54 & 6,08  & 41,72 & 9,08 \\
\midrule
\multirow{3}{*}{Vị trí kệ (shelf)}
  & Kệ 1 (Thấp) & 20 & 46,15 & 13,56 & 43,93 & 15,07 \\
  & Kệ 2 (Vừa)  & 20 & 33,98 & 13,53 & 30,58 & 15,71 \\
  & Kệ 3 (Cao)  & 36 & 45,22 & 12,97 & 40,80 & 15,48 \\
\bottomrule
\end{tabular}
\end{table}

\begin{figure}[H]
\centering
\includegraphics[width=\textwidth]{Hinh_3_GroupedBoxplot.png}
\caption{So sánh phân phối điểm đánh giá (rating) theo Nhà sản xuất và Vị trí kệ (N = 76).
Hình thoi trắng: trung bình nhóm.}
\label{fig:group}
\end{figure}

\begin{table}[H]
\centering
\caption{Hàm lượng đường và chất xơ (trên 1 oz) theo vị trí kệ}
\label{tab:duongke}
\begin{tabular}{l c r r r}
\toprule
Vị trí kệ & $n$ & Mean sugars & Median sugars & Mean fiber \\
\midrule
Kệ 1 (Thấp) & 20 & 5,20 & 3,00  & 1,72 \\
Kệ 2 (Vừa)  & 20 & 9,80 & 11,50 & 0,89 \\
Kệ 3 (Cao)  & 36 & 5,83 & 6,00  & 2,94 \\
\bottomrule
\end{tabular}
\end{table}

Quan sát Bảng~\ref{tab:nhom}, Bảng~\ref{tab:duongke} và Hình~\ref{fig:group}(a), (b) rút ra hai
nhận định trọng tâm (mới ở mức mô tả, chưa qua kiểm định).

\textbf{Chênh lệch rõ rệt theo biến mfr:} Toàn bộ hộp tứ phân vị của Nabisco nằm tách
biệt ở vùng điểm cao ($\xbar = 67{,}97$; $\xmed = 68{,}32$) với độ phân tán rất hẹp
($SD = 5{,}51$; $IQR = 6{,}24$), cao gần gấp đôi so với General Mills ở vị trí thấp nhất
($\xbar = 34{,}49$; $\xmed = 36{,}18$). Lưu ý Nabisco chỉ có $n = 6$. Nhóm Kelloggs có trung
bình bị kéo lên cao hơn trung vị ($44{,}04$ so với $40{,}56$) do điểm ngoại lai đỏ tại 93,70
điểm, trong khi Quaker Oats có độ biến động nội bộ lớn nhất ($SD = 16{,}80$; $IQR = 22{,}35$).

\textbf{Sự phân hóa theo Vị trí kệ (shelf):} Sản phẩm đặt tại Kệ 2 (Vừa) có rating
thấp hơn ($\xbar = 33{,}98$) so với Kệ 1 (Thấp, $\xbar = 46{,}15$) và Kệ 3 (Cao,
$\xbar = 45{,}22$). Bảng~\ref{tab:duongke} cho thấy Kệ 2 cũng là nơi trưng bày nhiều ngũ cốc
nhiều đường nhất (đường trung bình 9,80 g/oz so với 5,20 và 5,83 g/oz ở kệ 1 và kệ 3) và ít
chất xơ nhất. Khác biệt này có ý nghĩa thống kê hay không sẽ do ANOVA/Kruskal-Wallis kết
luận; theo mục~\ref{sec:rating}, nếu có, nó phản ánh thành phần dinh dưỡng của các sản phẩm
trên từng kệ.

\subsubsection{Ma trận tương quan và Biểu đồ phân tán tiêu biểu}
\label{sec:tuongquan}

\begin{figure}[H]
\centering
\includegraphics[width=\textwidth]{Hinh_4_Correlation_Scatter.png}
\caption{Ma trận tương quan Pearson và Biểu đồ phân tán (Scatter Plot) của điểm đánh giá (N = 76)}
\label{fig:corr}
\end{figure}

Vì hệ số Pearson nhạy với ngoại lai, trong khi nhóm đã quyết định giữ các ngoại lai
(mục 5.1.3), Bảng~\ref{tab:pearsonspearman} báo cáo song song hệ số Pearson, hệ số Spearman và
hệ số Pearson khi tạm bỏ 3 sản phẩm cám (100\% Bran, All-Bran, All-Bran with Extra Fiber).

\begin{table}[H]
\centering
\caption{Hệ số tương quan của rating với từng biến: Pearson, Spearman và Pearson khi bỏ 3 sản phẩm cám}
\label{tab:pearsonspearman}
\begin{tabular}{l r r r}
\toprule
Biến & Pearson ($N = 76$) & Spearman ($N = 76$) & Pearson bỏ cám ($N = 73$) \\
\midrule
sugars       & $-0{,}765$ & $-0{,}823$ & $-0{,}791$ \\
calories     & $-0{,}648$ & $-0{,}648$ & $-0{,}511$ \\
total\_carbs & $-0{,}585$ & $-0{,}491$ & $-0{,}414$ \\
fat          & $-0{,}397$ & $-0{,}431$ & $-0{,}406$ \\
sodium       & $-0{,}360$ & $-0{,}195$ & $-0{,}426$ \\
vitamins     & $-0{,}240$ & $-0{,}332$ & $-0{,}262$ \\
carbo        & $0{,}255$  & $0{,}395$  & $0{,}529$ \\
potass       & $0{,}524$  & $0{,}425$  & $0{,}325$ \\
protein      & $0{,}574$  & $0{,}653$  & $0{,}534$ \\
fiber        & $0{,}660$  & $0{,}584$  & $0{,}518$ \\
\bottomrule
\end{tabular}
\end{table}

Từ Hình~\ref{fig:corr}(a), (b) và Bảng~\ref{tab:pearsonspearman}, mối liên hệ giữa điểm đánh
giá rating và các thành phần dinh dưỡng thể hiện hai chiều hướng đối lập rõ rệt (phân mức theo
quy ước ở mục 3.5.1).

\textbf{Nhóm tương quan nghịch:} Chỉ có hàm lượng đường sugars đạt mức tương quan âm
\emph{mạnh} ($r = -0{,}765$; Spearman $-0{,}823$) và ổn định khi bỏ 3 sản phẩm cám. Năng
lượng calories ($r = -0{,}648$) và total\_carbs ($-0{,}585$) ở mức \emph{trung bình}; fat,
sodium và vitamins tương quan âm yếu. Trên Hình~\ref{fig:corr}(b), các điểm quan sát bám khá
sát đường hồi quy dốc xuống, cho thấy sản phẩm càng nhiều đường và năng lượng thì rating càng
thấp.

\textbf{Nhóm tương quan thuận:} Chất xơ fiber ($r = 0{,}660$), chất đạm protein
($0{,}574$) và kali potass ($0{,}524$) có tương quan dương ở mức \emph{trung bình} với rating;
carbo tương quan thuận yếu.

\textbf{Ảnh hưởng của cụm sản phẩm cám:} Tương quan của fiber, potass và calories với
rating giảm rõ khi bỏ 3 sản phẩm cám (potass từ 0,524 xuống 0,325; fiber từ 0,660 xuống
0,518), và hệ số Spearman của các biến này cũng thấp hơn Pearson. Như vậy một phần độ lớn
của các hệ số Pearson này đến từ 3 sản phẩm cám, trong khi tương quan của sugars vẫn mạnh
theo cả ba cách tính.

\textbf{Chiều tương quan không phải chiều tác động:} potass tương quan thuận với rating,
nhưng trong công thức tính rating (mục~\ref{sec:rating}) hệ số của potass lại âm
($-0{,}034$), vì potass luôn đi kèm fiber. Do đó ``kali tương quan thuận với rating'' là mô
tả đúng, nhưng không được diễn giải thành ``kali làm tăng rating'' ở phần hồi quy.

\textbf{Nguy cơ đa cộng tuyến:} Ma trận tương quan Hình~\ref{fig:corr}(a) ghi nhận cặp
fiber -- potass có hệ số tương quan rất mạnh ($r = 0{,}923$; Spearman 0,879; vẫn 0,876 khi bỏ
3 sản phẩm cám), cùng các cặp total\_carbs -- fiber ($-0{,}838$) và total\_carbs -- potass
($-0{,}804$). Tính trên 9 biến dinh dưỡng (dữ liệu 1 oz), VIF của carbo (13,5), sugars (11,7)
và fiber (10,5) đã vượt ngưỡng 10, potass (8,9) và calories (7,1) vượt ngưỡng 5. Đa cộng tuyến
vì vậy ở mức nặng và cần được xử lý (bỏ bớt hoặc gộp biến) trước khi chốt mô hình hồi quy.

% =====================================================================
\section{THỐNG KÊ SUY DIỄN}
\label{sec:suydien}
% =====================================================================
% Mục 6.1 và 6.2 (kiểm định một mẫu, hai mẫu) do thành viên khác phụ trách,
% sẽ được chèn vào trước mục này. Bộ đếm được đặt để mục dưới đây là 6.3.
\setcounter{subsection}{2}

\subsection{Phân tích phương sai một yếu tố (One-way ANOVA)}
\label{sec:anova}

\subsubsection{Mục tiêu}

\textbf{Biến phụ thuộc và yếu tố phân nhóm:} Biến phụ thuộc là điểm đánh giá
\texttt{rating}. Hai yếu tố phân nhóm được phân tích riêng rẽ: Nhà sản xuất \texttt{mfr}
($k = 6$, cỡ nhóm từ 6 đến 23) và Vị trí kệ \texttt{shelf} ($k = 3$, cỡ nhóm 20, 20, 36). Để
tìm nguyên nhân dinh dưỡng đứng sau, nhóm lặp lại cùng quy trình cho hàm lượng đường
\texttt{sugars} (g/oz), biến có tương quan mạnh nhất với rating ở Bảng~\ref{tab:pearsonspearman}.

\textbf{Mô hình (mục 3.4.1):} $Y_{ij} = \mu + \tau_i + \varepsilon_{ij}$, trong đó
$Y_{ij}$ là rating của sản phẩm thứ $j$ thuộc nhóm $i$, $\tau_i$ là hiệu ứng của nhóm $i$.

\textbf{Cặp giả thuyết (với mức ý nghĩa $\alpha = 0{,}05$):}
\begin{equation}
\begin{cases}
H_0: \mu_1 = \mu_2 = \dots = \mu_k\\
H_1: \text{Tồn tại ít nhất một cặp nhóm có trung bình khác nhau}
\end{cases}
\end{equation}
Với Kruskal-Wallis, $H_0$ được phát biểu là phân phối rating (hay hạng trung bình) như nhau
giữa các nhóm.

\textbf{So với Mục 6.2:} So sánh từng cặp cho 6 hãng cần tới 15 kiểm định $t$ riêng lẻ,
làm tăng xác suất kết luận sai. ANOVA kiểm định tất cả các nhóm trong một lần, sau đó các
kiểm định hậu nghiệm (Tukey HSD, Dunn) chỉ ra cặp khác biệt mà vẫn kiểm soát sai lầm loại I
toàn họ.

\subsubsection{Kết quả}

\textbf{Bước 1 -- Kiểm tra giả định (mục 3.4.2 -- 3.4.3).} Kết quả Shapiro-Wilk trên từng nhóm
được trình bày ở Bảng~\ref{tab:shapiro}, kiểm định Levene (biến thể Brown-Forsythe) ở
Bảng~\ref{tab:levene}.

\begin{table}[H]
\centering
\caption{Kiểm định Shapiro-Wilk cho rating trong từng nhóm}
\label{tab:shapiro}
\begin{tabular}{l l c c c}
\toprule
Yếu tố & Nhóm & $n_i$ & $W$ & $p$-value \\
\midrule
\multirow{6}{*}{mfr}
  & General Mills  & 22 & 0,9569 & 0,4291 \\
  & Kelloggs       & 23 & 0,8046 & $4{,}55\cdot10^{-4}$ \\
  & Nabisco        & 6  & 0,9568 & 0,7949 \\
  & Post           & 9  & 0,9013 & 0,2596 \\
  & Quaker Oats    & 8  & 0,9227 & 0,4523 \\
  & Ralston Purina & 8  & 0,9186 & 0,4186 \\
\midrule
\multirow{3}{*}{shelf}
  & Kệ 1 (Thấp) & 20 & 0,9116 & 0,0682 \\
  & Kệ 2 (Vừa)  & 20 & 0,8752 & 0,0145 \\
  & Kệ 3 (Cao)  & 36 & 0,8588 & $3{,}05\cdot10^{-4}$ \\
\bottomrule
\end{tabular}
\end{table}

\begin{table}[H]
\centering
\caption{Kiểm định Levene (Brown-Forsythe) về tính đồng nhất phương sai}
\label{tab:levene}
\begin{tabular}{l c c c c}
\toprule
Bài toán & $W$ & $df_1$ & $df_2$ & $p$-value \\
\midrule
rating theo mfr   & 1,426 & 5 & 70 & 0,2257 \\
rating theo shelf & 0,170 & 2 & 73 & 0,8441 \\
sugars theo mfr   & 0,513 & 5 & 70 & 0,7654 \\
sugars theo shelf & 1,133 & 2 & 73 & 0,3277 \\
\bottomrule
\end{tabular}
\end{table}

\textbf{Phân phối chuẩn:} Giả định chuẩn bị vi phạm ở Kelloggs ($p < 0{,}001$), Kệ 2
($p = 0{,}0145$) và Kệ 3 ($p < 0{,}001$). Đây đều là các nhóm chứa sản phẩm cám có rating rất
cao (All-Bran with Extra Fiber 93,70 điểm thuộc Kelloggs, kệ 3). Nabisco có $p = 0{,}7949$
nhưng chỉ có $n = 6$, nên kết quả này không đủ để khẳng định nhóm có phân phối chuẩn.

\textbf{Đồng nhất phương sai:} Cả bốn bài toán đều có $p$-value của Levene lớn hơn 0,05,
chưa có cơ sở bác bỏ giả thuyết phương sai bằng nhau.

\textbf{Hướng xử lý:} Do giả định chuẩn bị vi phạm ở một số nhóm và cỡ nhóm mất cân bằng
(6 đến 36), nhóm báo cáo song song ba kiểm định: ANOVA cổ điển, Welch's ANOVA (không cần
phương sai bằng nhau) và Kruskal-Wallis (không cần phân phối chuẩn). Kết luận chỉ được đưa ra
khi các kiểm định thống nhất; khi chúng khác nhau, nhóm ưu tiên Kruskal-Wallis và Dunn.

\textbf{Bước 2 -- Kiểm định chính (mục 3.4.4, 3.4.6).} Bảng~\ref{tab:anova_rating} trình bày
bảng ANOVA theo đúng cấu trúc Bảng~\ref{tab:anova}. Ngoài các đại lượng chuẩn, nhóm báo cáo thêm
tỷ lệ biến thiên do yếu tố nhóm $\eta^2 = SSB/SST$ để đo độ lớn của khác biệt.

\begin{table}[H]
\centering
\caption{Bảng phân tích phương sai một yếu tố cho rating}
\label{tab:anova_rating}
\begin{tabular}{l l r r r r c}
\toprule
Yếu tố & Nguồn biến thiên & SS & df & MS & $F$ & $p$-value \\
\midrule
\multirow{3}{*}{\makecell[l]{mfr\\($\eta^2 = 0{,}362$)}}
  & Giữa các nhóm & 5\,373,70  & 5  & 1\,074,74 & 7,942 & $5{,}68\cdot10^{-6}$ \\
  & Nội bộ nhóm   & 9\,472,67  & 70 & 135,32    &       & \\
  & Toàn phần     & 14\,846,37 & 75 &           &       & \\
\midrule
\multirow{3}{*}{\makecell[l]{shelf\\($\eta^2 = 0{,}134$)}}
  & Giữa các nhóm & 1\,984,31  & 2  & 992,15 & 5,631 & 0,0053 \\
  & Nội bộ nhóm   & 12\,862,06 & 73 & 176,19 &       & \\
  & Toàn phần     & 14\,846,37 & 75 &        &       & \\
\bottomrule
\end{tabular}
\end{table}

Bảng~\ref{tab:omnibus} tổng hợp ba kiểm định tổng thể cho cả bốn bài toán. Vì nhóm thực hiện 4
kiểm định Kruskal-Wallis cùng lúc, cột cuối hiệu chỉnh $p$-value theo phương pháp Holm
(mục 3.2.1).

\begin{table}[H]
\centering
\caption{Kiểm định tổng thể: ANOVA cổ điển, Welch's ANOVA và Kruskal-Wallis ($N = 76$)}
\label{tab:omnibus}
\small
\setlength{\tabcolsep}{3.5pt}
\begin{tabular}{l c c c c c c c c c}
\toprule
 & \multicolumn{2}{c}{ANOVA} & \multicolumn{3}{c}{Welch's ANOVA} & \multicolumn{4}{c}{Kruskal-Wallis} \\
\cmidrule(lr){2-3}\cmidrule(lr){4-6}\cmidrule(lr){7-10}
Bài toán & $F$ & $p$ & $F$ & $df_2$ & $p$ & $H_c$ & $df$ & $p$ & $p$ (Holm) \\
\midrule
rating theo mfr   & 7,942 & $5{,}68\cdot10^{-6}$ & 24,806 & 23,29 & $1{,}29\cdot10^{-8}$ & 22,20 & 5 & $4{,}80\cdot10^{-4}$ & 0,0014 \\
rating theo shelf & 5,631 & 0,0053 & 5,388 & 40,33 & 0,0084 & 13,32 & 2 & 0,0013 & 0,0026 \\
sugars theo mfr   & 2,425 & 0,0437 & 3,507 & 21,99 & 0,0176 & 10,60 & 5 & 0,0599 & 0,0599 \\
sugars theo shelf & 9,878 & $1{,}60\cdot10^{-4}$ & 8,709 & 36,25 & $8{,}15\cdot10^{-4}$ & 15,94 & 2 & $3{,}45\cdot10^{-4}$ & 0,0014 \\
\bottomrule
\end{tabular}
\end{table}

\textbf{Bước 3 -- Phân tích hậu nghiệm (mục 3.4.5 -- 3.4.6).} Vì các kiểm định tổng thể cho
rating đều bác bỏ $H_0$, nhóm so sánh từng cặp bằng Tukey HSD (khoảng tin cậy 95\% đồng thời) và
Dunn's test (hiệu chỉnh Bonferroni) trong Bảng~\ref{tab:posthoc}; khoảng tin cậy Tukey được biểu
diễn trên Hình~\ref{fig:tukey}.

\begin{table}[H]
\centering
\caption{So sánh bội hậu nghiệm cho rating: Tukey HSD và Dunn's test (Bonferroni)}
\label{tab:posthoc}
\small
\setlength{\tabcolsep}{3.5pt}
\begin{tabular}{l l r c c r c}
\toprule
 & & \multicolumn{3}{c}{Tukey HSD} & \multicolumn{2}{c}{Dunn (Bonferroni)} \\
\cmidrule(lr){3-5}\cmidrule(lr){6-7}
Yếu tố & Cặp so sánh & Chênh lệch & KTC 95\% & $p_{\text{adj}}$ & $z$ & $p_{\text{adj}}$ \\
\midrule
\multirow{15}{*}{mfr}
 & Kelloggs -- General Mills       & 9,55   & [$-0{,}61$; 19,72]   & 0,0775 & 2,376  & 0,2624 \\
 & Nabisco -- General Mills        & 33,48  & [17,78; 49,18]       & $4{,}18\cdot10^{-7}$ & 4,596 & $6{,}47\cdot10^{-5}$ \\
 & Post -- General Mills           & 7,22   & [$-6{,}27$; 20,71]   & 0,6217 & 1,630  & 1,0000 \\
 & Quaker Oats -- General Mills    & 8,43   & [$-5{,}64$; 22,50]   & 0,5009 & 1,718  & 1,0000 \\
 & Ralston Purina -- General Mills & 7,06   & [$-7{,}02$; 21,13]   & 0,6846 & 1,718  & 1,0000 \\
 & Nabisco -- Kelloggs             & 23,93  & [8,30; 39,56]        & $3{,}85\cdot10^{-4}$ & 3,072 & 0,0319 \\
 & Post -- Kelloggs                & $-2{,}33$  & [$-15{,}73$; 11,07]  & 0,9956 & $-0{,}162$ & 1,0000 \\
 & Quaker Oats -- Kelloggs         & $-1{,}12$  & [$-15{,}11$; 12,87]  & 0,9999 & 0,001  & 1,0000 \\
 & Ralston Purina -- Kelloggs      & $-2{,}50$  & [$-16{,}49$; 11,50]  & 0,9951 & 0,001  & 1,0000 \\
 & Post -- Nabisco                 & $-26{,}26$ & [$-44{,}23$; $-8{,}30$] & $7{,}91\cdot10^{-4}$ & $-2{,}792$ & 0,0785 \\
 & Quaker Oats -- Nabisco          & $-25{,}05$ & [$-43{,}46$; $-6{,}64$] & 0,0022 & $-2{,}606$ & 0,1373 \\
 & Ralston Purina -- Nabisco       & $-26{,}43$ & [$-44{,}83$; $-8{,}02$] & 0,0010 & $-2{,}606$ & 0,1373 \\
 & Quaker Oats -- Post             & 1,21   & [$-15{,}35$; 17,77]  & 0,9999 & 0,132  & 1,0000 \\
 & Ralston Purina -- Post          & $-0{,}16$  & [$-16{,}73$; 16,40]  & 1,0000 & 0,132  & 1,0000 \\
 & Ralston Purina -- Quaker Oats   & $-1{,}37$  & [$-18{,}42$; 15,67]  & 0,9999 & 0,000  & 1,0000 \\
\midrule
\multirow{3}{*}{shelf}
 & Kệ 2 (Vừa) -- Kệ 1 (Thấp) & $-12{,}17$ & [$-22{,}21$; $-2{,}12$] & 0,0136 & $-3{,}115$ & 0,0055 \\
 & Kệ 3 (Cao) -- Kệ 1 (Thấp) & $-0{,}93$  & [$-9{,}78$; 7,93]        & 0,9662 & $-0{,}198$ & 1,0000 \\
 & Kệ 3 (Cao) -- Kệ 2 (Vừa)  & 11,24      & [2,38; 20,10]            & 0,0092 & 3,334      & 0,0026 \\
\bottomrule
\end{tabular}
\end{table}

\begin{figure}[H]
\centering
\includegraphics[width=\textwidth]{Hinh_6_3_Tukey_HSD.png}
\caption{Khoảng tin cậy 95\% đồng thời Tukey HSD cho chênh lệch rating trung bình giữa các cặp
nhóm. Màu đỏ: khoảng không chứa 0 ($p_{\text{adj}} < 0{,}05$).}
\label{fig:tukey}
\end{figure}

Để tìm nguyên nhân dinh dưỡng, Bảng~\ref{tab:sugars_nhom} so sánh hàm lượng đường giữa các nhóm,
kèm $p$-value đã hiệu chỉnh của Dunn's test cho các cặp đáng chú ý.

\begin{table}[H]
\centering
\caption{Hàm lượng đường sugars (g/oz) theo nhóm và các cặp khác biệt (Dunn, Bonferroni)}
\label{tab:sugars_nhom}
\begin{tabular}{l l c r r r}
\toprule
Yếu tố & Nhóm & $n_i$ & Mean & SD & Median \\
\midrule
\multirow{6}{*}{mfr}
  & General Mills  & 22 & 7,56 & 3,63 & 8,00 \\
  & Kelloggs       & 23 & 6,95 & 4,14 & 6,92 \\
  & Nabisco        & 6  & 1,83 & 2,86 & 0,00 \\
  & Post           & 9  & 8,17 & 4,17 & 8,00 \\
  & Quaker Oats    & 8  & 6,25 & 4,43 & 6,50 \\
  & Ralston Purina & 8  & 6,13 & 3,56 & 5,50 \\
\midrule
\multirow{3}{*}{shelf}
  & Kệ 1 (Thấp) & 20 & 5,20 & 4,38 & 3,00 \\
  & Kệ 2 (Vừa)  & 20 & 9,80 & 3,92 & 11,50 \\
  & Kệ 3 (Cao)  & 36 & 5,83 & 3,03 & 6,00 \\
\bottomrule
\end{tabular}

\vspace{0.4em}
\begin{minipage}{0.8\textwidth}\small
Dunn's test (Bonferroni) cho sugars. Theo kệ: Kệ 2 -- Kệ 1 $p_{\text{adj}} = 9{,}29\cdot10^{-4}$;
Kệ 3 -- Kệ 2 $p_{\text{adj}} = 0{,}0016$; Kệ 3 -- Kệ 1 $p_{\text{adj}} = 1{,}0000$. Theo hãng:
Nabisco -- General Mills $p_{\text{adj}} = 0{,}0343$; Nabisco -- Post $0{,}0750$;
Nabisco -- Kelloggs $0{,}1294$; các cặp còn lại $p_{\text{adj}} > 0{,}6$.
\end{minipage}
\end{table}

\subsubsection{Nhận xét}

\textbf{Rating khác biệt theo nhà sản xuất:} Cả ba kiểm định đều bác bỏ $H_0$ ở mức
0,05: $F = 7{,}942 > F_{0,05;\,5,\,70} = 2{,}35$ ($p = 5{,}68\cdot10^{-6}$), Welch
$p = 1{,}29\cdot10^{-8}$ và Kruskal-Wallis $H_c = 22{,}20 > \chi^2_{0,05;\,5} = 11{,}07$
($p$ sau hiệu chỉnh Holm $= 0{,}0014$). Yếu tố hãng giải thích khoảng 36,2\% biến thiên của
rating ($\eta^2 = 0{,}362$).

\textbf{Nabisco là nguồn gốc của khác biệt:} Theo Tukey HSD, Nabisco có rating trung
bình cao hơn cả 5 hãng còn lại, chênh từ 23,93 điểm (so với Kelloggs) đến 33,48 điểm (so với
General Mills), và cả 5 khoảng tin cậy đều không chứa 0. Mọi cặp còn lại giữa 5 hãng không khác
biệt có ý nghĩa ($p_{\text{adj}} \ge 0{,}077$). Tuy nhiên, Dunn's test (vững hơn khi dữ liệu
không chuẩn) chỉ xác nhận Nabisco cao hơn General Mills ($p_{\text{adj}} < 0{,}001$) và
Kelloggs ($p_{\text{adj}} = 0{,}032$); ba cặp với Post, Quaker Oats và Ralston Purina có $p$
chưa hiệu chỉnh từ 0,005 đến 0,009 nhưng không còn ý nghĩa sau hiệu chỉnh Bonferroni. Lý do là
Nabisco chỉ có 6 sản phẩm và các hãng kia chỉ 8 -- 9 sản phẩm, nên lực kiểm định thấp. Vì vậy,
nhận định ``Nabisco tách biệt hoàn toàn khỏi các hãng còn lại'' chỉ được ủng hộ đầy đủ bởi
Tukey HSD; với cách kiểm định thận trọng hơn, bằng chứng chắc chắn chỉ có cho hai hãng lớn
nhất là General Mills và Kelloggs.

\textbf{Rating khác biệt theo vị trí kệ:} Cả ba kiểm định đều bác bỏ $H_0$:
$F = 5{,}631 > F_{0,05;\,2,\,73} = 3{,}12$ ($p = 0{,}0053$), Welch $p = 0{,}0084$, Kruskal-Wallis
$H_c = 13{,}32 > \chi^2_{0,05;\,2} = 5{,}99$ ($p$ Holm $= 0{,}0026$). Yếu tố kệ giải thích
13,4\% biến thiên của rating. Cả Tukey HSD và Dunn đều cho cùng kết luận: Kệ 2 (Vừa) thấp hơn
Kệ 1 khoảng 12,17 điểm (KTC 95\%: 2,12 đến 22,21 điểm) và thấp hơn Kệ 3 khoảng 11,24 điểm
(KTC 95\%: 2,38 đến 20,10 điểm), còn Kệ 1 và Kệ 3 không khác biệt ($p_{\text{adj}} = 0{,}966$).

\textbf{Nguyên nhân dinh dưỡng -- đường:} Hàm lượng đường khác biệt có ý nghĩa theo kệ
(Kruskal-Wallis $p$ Holm $= 0{,}0014$): sản phẩm Kệ 2 chứa trung bình 9,80 g/oz (trung vị
11,50 g/oz), cao hơn có ý nghĩa so với Kệ 1 (5,20 g/oz) và Kệ 3 (5,83 g/oz). Đây đúng là hai
cặp kệ có rating khác biệt, phù hợp với tương quan âm mạnh giữa sugars và rating
($r = -0{,}765$). Theo nhà sản xuất, Nabisco có lượng đường thấp nhất (trung bình 1,83 g/oz,
trung vị 0 g/oz do các dòng lúa mì nguyên hạt không đường), nhưng bằng chứng thống kê yếu hơn:
ANOVA ($p = 0{,}044$) và Welch ($p = 0{,}018$) bác bỏ $H_0$, còn Kruskal-Wallis thì không
($p = 0{,}060$). Vì nhóm Nabisco vi phạm giả định chuẩn với sugars (Shapiro-Wilk
$p = 0{,}0036$), nhóm ưu tiên Kruskal-Wallis và kết luận khác biệt về đường giữa các hãng
\emph{chưa đủ} ý nghĩa thống kê; riêng cặp Nabisco -- General Mills có ý nghĩa theo Dunn
($p_{\text{adj}} = 0{,}034$).

\textbf{Cách hiểu đúng kết quả:} Theo mục~\ref{sec:rating}, rating được tính từ thành
phần dinh dưỡng. Do đó, rating khác biệt giữa các hãng và các kệ không có nghĩa thương hiệu hay
vị trí kệ trực tiếp làm thay đổi rating; nó phản ánh việc các hãng và các kệ chứa những sản
phẩm có thành phần dinh dưỡng khác nhau (ví dụ Kệ 2 tập trung ngũ cốc nhiều đường). Ngoài ra,
mẫu không được chọn ngẫu nhiên (mục 1.1), nên kết luận chỉ áp dụng cho nhóm sản phẩm tương tự.

\subsubsection{Kết luận}

Với mức ý nghĩa 5\%, có đủ bằng chứng thống kê để kết luận rằng điểm đánh giá trung bình của
ngũ cốc \textbf{khác nhau giữa các nhà sản xuất} và \textbf{khác nhau giữa các vị trí kệ}; kết
luận này thống nhất qua cả ANOVA cổ điển, Welch's ANOVA và Kruskal-Wallis, kể cả sau khi hiệu
chỉnh cho nhiều kiểm định. Theo nhà sản xuất, khác biệt đến từ Nabisco, hãng có rating cao hơn
rõ rệt so với General Mills và Kelloggs (các hãng còn lại tương đương nhau). Theo vị trí kệ,
Kệ 2 (Vừa) có rating thấp hơn có ý nghĩa so với cả Kệ 1 và Kệ 3, đồng thời cũng là kệ có lượng
đường cao hơn có ý nghĩa. Mức độ liên hệ của yếu tố hãng ($\eta^2 = 36{,}2\%$) lớn hơn yếu tố kệ
($\eta^2 = 13{,}4\%$).

Tuy vậy, ANOVA chỉ cho biết \emph{nhóm nào} có rating cao hay thấp, và kết quả về đường cho
thấy khác biệt giữa các nhóm thực chất là khác biệt về thành phần dinh dưỡng. Câu hỏi tiếp theo
là: \emph{mỗi thành phần dinh dưỡng cụ thể (đường, chất xơ, chất đạm\ldots) liên hệ với rating
như thế nào khi xét đồng thời}. Câu hỏi này được trả lời bằng mô hình hồi quy tuyến tính bội ở
Mục 6.4.

% =====================================================================
\begin{thebibliography}{9}
\addcontentsline{toc}{section}{TÀI LIỆU THAM KHẢO}
\bibitem{kaggle} Chris Crawford. \textit{80 Cereals: Nutrition data on 80 cereal products}.
Kaggle. \url{https://www.kaggle.com/datasets/crawford/80-cereals}
\bibitem{r} R Core Team. \textit{R: A Language and Environment for Statistical Computing}.
R Foundation for Statistical Computing, Vienna, Austria. \url{https://www.R-project.org/}
\bibitem{e1071} D. Meyer và cộng sự. \textit{e1071: Misc Functions of the Department of
Statistics, Probability Theory Group, TU Wien}. R package.
\bibitem{ggplot2} H. Wickham. \textit{ggplot2: Elegant Graphics for Data Analysis}.
Springer-Verlag New York, 2016.
\end{thebibliography}

\end{document}
